\ifx\danslelivre\undefined
\RequirePackage{luatex85}
\documentclass[10pt]{../fichiers-configuration-latex/smfart}
\usepackage{polyglossia}
%\usepackage[francais]{babel}
\usepackage{fontspec}
\usepackage[math-style=ISO,bold-style=upright,sans-style=upright]{unicode-math}
%\usepackage[mathletters]{ucs}
%\usepackage[utf8]{inputenc}
%\usepackage[T1]{fontenc}
\usepackage[unicode,colorlinks=true,linkcolor=blue,citecolor=Sepia,urlcolor=magenta]{hyperref}
\setmainlanguage{french}
\setmainfont[Ligatures=TeX]{Linux Libertine O}
% variante possible : TeXGyrePagella ?
\setmathfont{XITS Math}

% maths
\usepackage{amsmath}
%\usepackage{amssymb}
%\usepackage{mathrsfs}
%\usepackage{stmaryrd}
%\usepackage{euscript}

% couleurs
\usepackage[usenames,dvipsnames]{xcolor}

% graphiques
\usepackage{graphicx}
\usepackage[all]{xy}
% Xy-pic 3.8.9's PDF driver assumes that \fontdimen8\textfont3 is a
% legacy TeX rule width.  Under LuaTeX with unicode-math/XITS Math it is
% instead 6.62pt, which turns every arrow shaft into a thick black bar.
% Override only xypdf's private cached width; do not alter the math font.
\makeatletter
\AtBeginDocument{%
  \@ifundefined{xP@preclw}{}{%
    \edef\xP@preclw{0.4pt}%
    \edef\xP@lw{\xP@dim\xP@preclw}%
    \PackageInfo{sga4-xypdf-compat}{Xy-pic PDF line width set to \xP@preclw}%
  }%
}
\makeatother
%\usepackage{tikz}
%\usetikzlibrary{matrix,arrows,calc}
%\tikzset{column sep=3em,row sep=2em,text height=1.5ex, text depth=0.25ex,>=angle 90}
%\tikzset{description/.style={fill=white,inner sep=1pt}}

% hyperliens, 'etiquettes
\usepackage{xr-hyper}
\usepackage{../fichiers-configuration-latex/smfthm}
%\ifx\cacheretiquettes\undefined
%\usepackage[inline]{showlabels} \else \fi % pour rendre visibles les 'etiquettes
%\usepackage{../fichiers-configuration-latex/makeidx} % pour faire un index
%\usepackage[pagebackref,unicode]{hyperref}  % pour voir o`u sont cit'ees les r'ef'erences
\hypersetup{unicode,colorlinks=true,linkcolor=blue,citecolor=Sepia,urlcolor=magenta,
  hypertexnames=false}
%\hypersetup{breaklinks=true} %permet le retour `a la ligne dans les liens trop longs. Incompatible avec pdflatex
%\renewcommand*{\backref}[1]{{\small [cit\'e page(s)~#1.]}}

\usepackage{xspace}
\usepackage{pifont} % pour expose XVII

\renewcommand{\thefootnote}{\roman{footnote}} % pour mettre les notes en bas de page sous la forme (ii) et non (2), qui ressemble trop (?) a un exposant.

\newcounter{pageoriginal}
\newcommand{\pageoriginale}{\refstepcounter{pageoriginal}\marginpar{\footnotesize\textbf{\thepageoriginal}}\xspace}
\let\pageoriginaled\pageoriginale

\usepackage{ifthen}
\newboolean{original}
\def\sisi{\ifthenelse{\boolean{original}}}
\setboolean{original}{false}

\newcounter{cnote}
\def\nde#1{\ifthenelse{\boolean{original}}{}{\stepcounter{cnote}\def\thefootnote{\arabic{cnote}}\protect\footnote[\value{cnote}]{N.D.E.: #1}}}
\def\ndemark{\ifthenelse{\boolean{original}}{}{\stepcounter{cnote}\def\thefootnote{\arabic{cnote}}\protect\footnotemark[\value{cnote}]}}
\def\ndetext#1{\ifthenelse{\boolean{original}}{}{\def\thefootnote{\arabic{cnote}}\protect\footnotetext[\value{cnote}]{N.D.E.: #1}}}

\setmainlanguage{french}
\setotherlanguage{english}

% unicode-math deliberately disables the legacy \Arrowvert glyph.  The
% source uses it once as a rotated equality marker; \Vert has the intended
% appearance and keeps the historical diagram buildable with LuaLaTeX.
\renewcommand{\Arrowvert}{\Vert}


%%
%% Pour afficher le numéro de version Git (exige que LuaTeX soit appelé avec --shell-escape)
%%

\catcode`\%=12
\directlua{
function command_output(cmd)
  local output_pipe = io.popen(cmd)
  local line_iter = output_pipe and output_pipe:lines()
  local line = line_iter and line_iter()
  if output_pipe then output_pipe:close() end
  return line
end
local git_command1 = [[git log --pretty=format:"%h %ai" -1]]
local git_command2 = [[if git diff --quiet HEAD ; then echo clean ; else echo dirty ; fi]]
function git_status()
  local status = command_output(git_command1)
  if status then
    tex.tprint({luatexbase.catcodetables.CatcodeTableLaTeX, string.char(0x5c).."texttt{"}, {-2,status}, {luatexbase.catcodetables.CatcodeTableLaTeX, "}"})
    if not (command_output(git_command2) == "clean") then tex.sprint(-2, " (modifications non validées)") end
  else
    tex.sprint(-2, "?")
  end
end
}
\catcode`\%=14
\def\showgitstatus{\directlua{git_status()}}


% macros pour LaTeX
\renewcommand{\mathsf}{}

% Police différente pour U+2113 SCRIPT SMALL L
\setmathfont[range="2113]{latinmodern-math.otf}

% Give every explicit SGA label its own destination and preserve the printed
% number supplied as the optional argument.  The former definition inherited
% whichever counter anchor happened to precede it and ignored #1 entirely.
\makeatletter
\def\etiquette[#1]#2{%
  \phantomsection
  \protected@edef\@currentlabel{#1}%
  \label{#2}%
}
\makeatother
\def\eqhyperref[#1]#2{(\hyperref[#1]{#2})}
\def\pagehyperref[#1]#2{(page \hyperref[#1]{#2})}
\def\hyperrefext[#1][#2]#3{\hyperref[#2]{#1 #3}}
\def\eqhyperrefext[#1][#2]#3{\hyperref[#2]{#1 #3}}     %cf. exposes 17
\def\neweqhyperrefext[#1][#2]#3{\hyperref[#2]{#1 #3}}  %cf. exposes 17

\def\shat#1{\protect\widehat{#1}}
\def\stilde#1{\protect\widetilde{#1}}
%%%%%%%%%%%%%%%%%%


\def\etcn{etc.\xspace}
\def\etco{etc...\xspace}
\def\etc{\sisi{\etco}{\etcn}}
\def\ldotsn{\sisi{...}{\etcn}}

\def\cf{{cf.}\kern.3em}
\def\Cf{{Cf.}\kern.3em}
\def\ie{{i.e.}\kern.3em}
\def\iev{{i.e.},\xspace}
\def\loccito{loc.\kern2pt cit.\xspace}
\def\loccitn{\textit{loc.\kern2pt cit.}\xspace}
\def\loccit{\sisi{\loccito}{\loccitn}}
\def\EGA{\textup{EGA}\kern.2em}
\def\SGA{\textup{SGA}\kern.3em}

\def\quoi{}

\def\quav{\quad,\quad}
\def\quaet{\quad \text{et}\quad}
\def\fibreproduct#1#2#3{#1 ×_{#3} #2}
\let\fprod\fibreproduct
\def\ccheck#1#2{\mbox{${\check{\check{\rm #1}}}^{\lower2pt\hbox{$\scriptstyle #2$}}$}}
\def\sheaf#1{\protect\mathcal{#1}}
\def\vieuxsheaf#1{{#1}} % cf. problème avec sheaf{a\Phi_i} par exemple...
\DeclareMathOperator{\Hom}{Hom}
\DeclareMathOperator{\can}{can}
\DeclareMathOperator{\Ind}{Ind}
\DeclareMathOperator{\ind}{ind}
%\DeclareMathOperator{\Hhomtop}{\mathcal{H}\!\mathit{omtop}}
\DeclareMathOperator{\Hhomtop}{\mathbf{Homtop}}
\DeclareMathOperator{\Point}{Point}
\DeclareMathOperator{\Pro}{Pro}
\DeclareMathOperator{\Top}{Top}
\DeclareMathOperator{\Mod}{Mod}
\DeclareMathOperator{\Cart}{Cart}
%\DeclareMathOperator{\Mmorsite}{\mathcal{M}\!\mathit{orsite}}
\DeclareMathOperator{\Mmorsite}{\mathbf{Morsite}}
\DeclareMathOperator{\loc}{loc}
\DeclareMathOperator{\Loc}{Loc}
\DeclareMathOperator{\Pic}{Pic}
\DeclareMathOperator{\Et}{Et}
\DeclareMathOperator{\card}{card}
\DeclareMathOperator{\Ens}{Ens}
\DeclareMathOperator{\h}{\underline{H}}
\DeclareMathOperator{\Fonct}{Fonct}
\DeclareMathOperator{\Ker}{Ker}
\DeclareMathOperator{\but}{but}
\DeclareMathOperator{\Ob}{Ob}
\DeclareMathOperator{\ob}{ob}
\DeclareMathOperator{\Lib}{Lib}
\DeclareMathOperator{\esj}{esj}
\DeclareMathOperator{\id}{id}
\DeclareMathOperator{\Id}{Id}
\DeclareMathOperator{\ad}{ad}
\DeclareMathOperator{\Iim}{Im}
\DeclareMathOperator{\pr}{pr}
\DeclareMathOperator{\Cat}{Cat}
\DeclareMathOperator{\car}{car}
\DeclareMathOperator{\SysInd}{SysInd}
\DeclareMathOperator{\Sup}{Sup}
\DeclareMathOperator{\Coker}{Coker}
\DeclareMathOperator{\Filt}{Filt}
\DeclareMathOperator{\acc}{acc}
\DeclareMathOperator{\sur}{sur}
\DeclareMathOperator{\cat}{cat}
\DeclareMathOperator{\Ord}{Ord}
\DeclareMathOperator{\Gr}{Gr}
\DeclareMathOperator{\Cov}{Cov}
\DeclareMathOperator{\Caf}{Caf}
\DeclareMathOperator{\coker}{coker}
\DeclareMathOperator{\coim}{coim}
\DeclareMathOperator{\im}{im}
\DeclareMathOperator{\Ouv}{Ouv}
\DeclareMathOperator{\diag}{diag}
\DeclareMathOperator{\Esp}{Esp}
\DeclareMathOperator{\TOP}{TOP}
\DeclareMathOperator{\Unif}{Unif}
\DeclareMathOperator{\Iind}{\mathbf{Ind}}
\DeclareMathOperator{\Hhomcart}{\mathbf{Homcart}}
\DeclareMathOperator{\Hhomtopess}{\mathbf{Homtopess}}
\DeclareMathOperator{\Ppoints}{\mathbf{Points}}
\DeclareMathOperator{\Ppoint}{\mathbf{Point}}
\DeclareMathOperator{\Ttors}{\mathbf{Tors}}
\DeclareMathOperator{\Ssex}{\mathbf{Sex}}
\DeclareMathOperator{\Ffib}{\mathbf{Fib}}
\DeclareMathOperator{\Vvois}{\mathbf{Vois}}
\DeclareMathOperator{\Ppointess}{\mathbf{Pointess}}
\DeclareMathOperator{\scrHom}{\mathbf{Hom}}
\DeclareMathOperator{\sheafPic}{\mathbf{Pic}}
\DeclareMathOperator{\sheafAut}{\mathbf{Aut}}
\DeclareMathOperator{\Hhomtoplocan}{\mathbf{Homtoplocan}}
\DeclareMathOperator{\Hhomtopan}{\mathbf{Homtopan}}
\DeclareMathOperator{\Pptgeom}{\mathbf{Ptgéom}}
\DeclareMathOperator{\Morsite}{Morsite}
\DeclareMathOperator{\Fib}{Fib}
\DeclareMathOperator{\Fais}{Fais}
\DeclareMathOperator{\kar}{kar}
\DeclareMathOperator{\Groupes}{Groupes}
\DeclareMathOperator{\Restr}{Restr}
\DeclareMathOperator{\Isom}{Isom}
\DeclareMathOperator{\Prol}{Prol}
\DeclareMathOperator{\Res}{Res}
\DeclareMathOperator{\Tors}{Tors}
\DeclareMathOperator{\esp}{esp}
\DeclareMathOperator{\es}{es}
\DeclareMathOperator{\od}{od}
\DeclareMathOperator{\gr}{gr}
\DeclareMathOperator{\sob}{sob}
\DeclareMathOperator{\dfn}{dfn}
\DeclareMathOperator{\Sch}{Sch}
\DeclareMathOperator{\ouv}{ouv}
\DeclareMathOperator{\cons}{cons}
\DeclareMathOperator{\Nerf}{Nerf}
\DeclareMathOperator{\Points}{Points}
\DeclareMathOperator{\Espsob}{Espsob}
\DeclareMathOperator{\Homtop}{Homtop}
\DeclareMathOperator{\Pt}{Pt}
\DeclareMathOperator{\Quot}{Quot}
\DeclareMathOperator{\glob}{glob}
\DeclareMathOperator{\Homfib}{Homfib}
\DeclareMathOperator{\Qouv}{Qouv}
\DeclareMathOperator{\Quotouv}{Quotouv}
\DeclareMathOperator{\ens}{ens}
\DeclareMathOperator{\ab}{ab}
\DeclareMathOperator{\an}{an}
\DeclareMathOperator{\fs}{fs}
\DeclareMathOperator{\red}{red}
\DeclareMathOperator{\Ab}{Ab}
\DeclareMathOperator{\cosq}{cosq}
\DeclareMathOperator{\Ann}{Ann}
\DeclareMathOperator{\Zar}{Zar}


%%%%%%%%%%%%
%%%%% SGA4-2
\DeclareMathOperator{\Sex}{Sex}
\DeclareMathOperator{\vois}{vois}
\DeclareMathOperator{\Vois}{Vois}
\DeclareMathOperator{\Ext}{Ext}
\DeclareMathOperator{\ext}{ext}
\DeclareMathOperator{\Fer}{Fer}
\DeclareMathOperator{\Add}{Add}
\DeclareMathOperator{\Cocart}{Cocart}
\DeclareMathOperator{\cosk}{cosk}
\DeclareMathOperator{\plat}{plat}

\DeclareMathOperator{\faisc}{faisc}
\DeclareMathOperator{\sk}{sk}
\DeclareMathOperator{\ch}{ch}

\def\HH{\mathbf{H}}
\def\qc{\mathrm{qc}}
\def\et{\mathrm{et}}
\def\coh{\mathrm{coh}}

\def\Hvr#1{%
\begin{array}[b]{l}
{\kern2pt\scriptstyle r}\\[-8pt]
{\kern1pt\scriptstyle \vee}\\[-5pt]
H^{#1}
\end{array}}

\def\Hvv#1{%
\begin{array}[b]{l}
\kern2pt\rotatebox{90}{${\scriptscriptstyle\ll}$}\\[-6pt]
H^{#1}
\end{array}}

\def\Hvi#1{%
\begin{array}[b]{l}
{\kern0.2pt\scriptscriptstyle \infty}\\[-8pt]
{\kern1pt\scriptstyle \vee}\\[-5pt]
H^{#1}
\end{array}}

\def\Hv#1{%
\begin{array}[b]{l}
{\kern1pt\scriptstyle \vee}\\[-5pt]
H^{#1}
\end{array}}

\DeclareMathOperator{\cart}{cart}
\DeclareMathOperator{\Ffaisc}{\mathbf{Faisc}}


\def\Ev#1{%
\begin{array}[b]{l}
{\kern1pt\scriptstyle \vee}\\[-5pt]
E'_{#1}
\end{array}}

\DeclareMathOperator{\Cont}{Cont}
\DeclareMathOperator{\Ccont}{\mathscr{C}\!\mathit{ont}}
\DeclareMathOperator{\Site}{Site}
\newcommand\MM{\mathbf{M}}
\def\limtop#1{{\displaystyle{\varprojlim_{#1}\!\mbox{\rm top }}}}
\def\Limtop{\mathop{\varprojlim\!\mbox{\rm top}}}
\DeclareMathOperator{\spec}{spec}
\def\ss{\scriptstyle{\mathscr{S}}}
\DeclareMathOperator{\Gal}{Gal}
\def\smallm{\scriptstyle{\mathscr{M}}}
\DeclareMathOperator{\End}{End}
\DeclareMathOperator{\ass}{ass}
\DeclareMathOperator{\zar}{zar}
\DeclareMathOperator{\Ogg}{Ogg}
\DeclareMathOperator{\res}{res}
\DeclareMathOperator{\tr}{tr}
\DeclareMathOperator{\Supp}{Supp}
\DeclareMathOperator{\codim}{codim}
\DeclareMathOperator{\rk}{rk}
\DeclareMathOperator{\cur}{cur}
\DeclareMathOperator{\cdqc}{cdqc}
\newcommand\DD{\mathbf{D}}
\DeclareMathOperator{\cl}{cl}
\def\JJ{\mathbf{J}}
\def\FF{\mathbf{F}}
\def\BB{\mathbf{B}}
\DeclareMathOperator{\Alget}{Alget}
\DeclareMathOperator{\Ass}{Ass}
\DeclareMathOperator{\supp}{supp}
\DeclareMathOperator{\rad}{rad}
\DeclareMathOperator{\prof}{prof}
\DeclareMathOperator{\cint}{int}
\DeclareMathOperator{\sic}{sic}
\DeclareMathOperator{\Tor}{Tor}
\DeclareMathOperator{\tors}{tors}
\DeclareMathOperator{\Sym}{Sym}
\DeclareMathOperator{\parf}{parf}
\DeclareMathOperator{\tor}{tor}
\DeclareMathOperator{\defn}{defn}
\DeclareMathOperator{\Div}{Div}
\DeclareMathOperator{\Pol}{Pol}
\DeclareMathOperator{\Fex}{Fex}
\DeclareMathOperator{\Lie}{Lie}
\DeclareMathOperator{\tot}{tot}
\DeclareMathOperator{\Aut}{Aut}
\DeclareMathOperator{\HOM}{HOM}
\DeclareMathOperator{\pch}{pch}
\DeclareMathOperator{\Ch}{Ch}
\DeclareMathOperator{\EXT}{EXT}
\DeclareMathOperator{\PIC}{PIC}
\DeclareMathOperator{\Rf}{Rf}
\DeclareMathOperator{\inj}{inj}
\newcommand\VV{\mathbf{V}}

\def\Xv{%
\begin{array}[b]{l}
{\kern1pt\scriptstyle \vee}\\[-5pt]
X
\end{array}}

%%%%%%%
\newcommand\TT{\mathbf{T}}
\newcommand\XX{\mathbf{X}}
\newcommand\EE{\mathbf{E}}
%%%%%%

\def\qedsymbol{\textup{C.Q.F.D.}}
\let\CQFD\qed
\let\cqfd\qed

%%%%%\bigdot
\font\bigsymb=cmsy10 at 4pt
\def\bigdot{{\kern1.2pt\raise 1.5pt\hbox{\bigsymb\char15}}}
%%%%%

\def\Lright{\mathop{\vcenter{\hbox{\hbox{$\longrightarrow$}\lower8pt \hbox{\kern-16pt\hbox{$\longrightarrow$}}}}}}

\def\Aa{{\mathcal A}}
\def\Bb{{\mathcal B}}
\def\Cc{{\mathcal C}}
\def\Dd{{\mathcal D}}
\def\Ee{{\mathcal E}}
\def\Ff{{\mathcal F}}
\def\Gg{{\mathcal G}}
\def\Hh{{\mathcal H}}
\def\Ii{{\mathcal I}}
\def\Jj{{\mathcal J}}
\def\Kk{{\mathcal K}}
\def\Ll{{\mathcal L}}
\def\Mm{{\mathcal M}}
\def\Nn{{\mathcal N}}
\def\Oo{{\mathcal O}}
\def\Pp{{\mathcal P}}
\def\Qq{{\mathcal Q}}
\def\Rr{{\mathcal R}}
\def\Ss{{\mathcal S}}
\def\Tt{{\mathcal T}}
\def\Uu{{\mathcal U}}
\def\Vv{{\mathcal V}}
\def\Ww{{\mathcal W}}
\def\Xx{{\mathcal X}}
\def\Yy{{\mathcal Y}}
\def\Zz{{\mathcal Z}}

\newcommand{\CC}{\mathbf{C}}
\newcommand{\GG}{\mathbf{G}}
\newcommand{\KK}{\mathbf{K}}
\newcommand{\LL}{\mathbf{L}}
\newcommand{\NN}{\mathbf{N}}
\newcommand{\PP}{\mathbf{P}}
\newcommand\QQ{\mathbf{Q}}
\newcommand\RR{\mathbf{R}}
\newcommand\ZZ{\mathbf{Z}}

\def\to{\mathchoice{\longrightarrow}{\rightarrow}{\rightarrow}{\rightarrow}}
\def\from{\mathchoice{\longleftarrow}{\leftarrow}{\leftarrow}{\leftarrow}}
\def\mto{\mathchoice{\longmapsto}{\mapsto}{\mapsto}{\mapsto}}
\def\mfrom{\mathrel{\from\joinrel\mapstochar\,}}
\def\hto{\mathrel{\lhook\joinrel\to}}
\def\hfrom{\mathrel{\from\joinrel\rhook}}

\def\implies{\mathchoice{\Longrightarrow}{\Rightarrow}{\Rightarrow}{\Rightarrow}}
\def\implied{\mathchoice{\Longleftarrow}{\Leftarrow}{\Leftarrow}{\Leftarrow}}
\def\iff{\mathchoice{\Longleftrightarrow}{\Leftrightarrow}{\Leftrightarrow}{\Leftrightarrow}}

\def\To#1{\mathchoice{\xrightarrow{\labelstyle~~#1~}}{\stackrel{#1}{\longrightarrow}}{}{}}
\def\Mto#1{\mathrel{\mapstochar\To{#1}}}
\def\From#1{\mathchoice{\xleftarrow{\labelstyle~~#1~}}{\stackrel{#1}{\longleftarrow}}{}{}}

\newcommand{\isomto}{\overset{\sim}{\longrightarrow}}
\newcommand{\isomfrom}{\overset{\sim}{\longleftarrow}}

\newcommand{\IFF}{\unskip~$\iff$~\ignorespaces}
\newcommand{\THEN}{\unskip~$\implies$~\ignorespaces}

\let\leq\leqslant
\let\geq\geqslant
\let\le\leqslant
\let\ge\geqslant

\def\cf{{cf.}\kern.3em}
\def\Cf{{Cf.}\kern.3em}
\def\ie{{i.e.}\kern.3em}
\def\iev{{i.e.},\xspace}
%\def\loccito{loc.\kern2pt cit.\xspace}
%\def\loccitn{\textit{loc.\kern2pt cit.}\xspace}
%\def\loccit{\sisi{\loccito}{\loccitn}}
\def\EGA{\textup{EGA}\kern.2em}
\def\SGA{\textup{SGA}\kern.3em}

\let\oldS\S
\def\S{\oldS\kern .15em}
\def\ptbl{.\kern .2em}
\def\resp{resp.\kern.3em}
\def\Exp{Exp.\kern.3em}

\let\bar\overline
\let\hat\widehat
\let\tilde\widetilde
\let\dpl\displaystyle
\let\ldots\dots
\let\emptyset\varnothing
\let\skpt\skippointrait
\let\epsilon\varepsilon

\def\sfrac#1#2{{#1}/{#2}}

\DeclareMathOperator{\cd}{cd}
%\DeclareMathOperator{\Hhom}{\mathcal{H}\!\mathit{om}}
\DeclareMathOperator{\Hhom}{\mathbf{Hom}}
%\DeclareMathOperator{\Eext}{\mathcal{E}\!\mathit{xt}}
\DeclareMathOperator{\Eext}{\mathbf{Ext}}
\DeclareMathOperator{\Spec}{Spec}
\DeclareMathOperator{\Tr}{Tr}

\let\emptyset=\varnothing

%%%%%%%%%%%%%%%
%%% expose XVII
%%%%%%%%%%%%%%%

\makeatletter
\let\old@optexpose\optexpose
\newcommand{\empty@optexpose}[1]{}
\newcommand{\paren@optexpose}[1]{\old@optexpose{#1}(}
\newcommand{\refnoexpose}[1]{\let\optexpose\empty@optexpose\ref{#1}\let\optexpose\old@optexpose}%\refnoexpose{X.cor5.2} engendre "5.2" au lieu de "X 5.2"
\newcommand{\neweqref}[1]{\let\optexpose\paren@optexpose\textup{\ref{#1})}\let\optexpose\old@optexpose}%\neweqref{XII.eq4.1} engendre "XII (4.1)" au lieu de "(XII 4.1)"
\makeatother
\newcommand{\refvb}[2][VI${}_{\text{B}}$ ]{\sisi{#1\refnoexpose{#2}}{\ref{#2}}}

\newcommand{\leftexp}[2]{{\vphantom{#2}}^{#1}{#2}}
\newcommand{\AAA}{\mathbf{A}}
\newcommand{\Sf}{\mathfrak{S}}
\DeclareMathOperator*{\quotprojlim}{\text{``$\varprojlim$''}}
\DeclareMathOperator*{\quotinjlim}{\text{``$\varinjlim$''}}
\DeclareMathOperator{\Fl}{F\sisi{\ell}{l}}
\DeclareMathOperator{\Pl}{P\sisi{\ell}{l}}
%\DeclareMathOperator{\cl}{c\sisi{\ell}{l}}
\DeclareMathOperator{\tdf}{tdf}
\DeclareMathOperator{\rg}{rg}
\newcommand{\iInd}{\sisi{I}{i}nd}
\newcommand{\iind}{ind}
\newcommand{\ppro}{pro}
\newcommand{\tTor}{\sisi{T}{t}or}
\newcommand{\ttor}{tor}
\newcommand{\chap}{Chap.}
\newcommand{\Ffex}{\Ff\mathit{ex}}
\newcommand{\sheafk}{\underline{k}}
\newcommand{\sheafn}{\underline{n}}
\newcommand{\sheafGamma}{\underline{\Gamma}}

\sisi{
\DeclareMathOperator{\simarrow}{\rule{.15mm}{1.8mm}\kern -3.4pt\rightsquigarrow}
\newcommand{\sheafell}{\underline{\ell}}
\newcommand{\sheafZZ}{\underline{\ZZ}}
}{
\newcommand{\nderef}[1]{la N.D.E.~(\ref{#1}) page~\pageref{#1}}
}

%%%%%%%%%%%%%%%%%
%%%%% expose XVII
%%%%%%%%%%%%%%%%%



\externaldocument{../01/01.tex}
\externaldocument{../02/02.tex}
\externaldocument{../03/03.tex}
\externaldocument{../04/04.tex}
\externaldocument{../05/05.tex}
\externaldocument{../05bis/05bis.tex}
\externaldocument{../06/06.tex}
\externaldocument{../07/07.tex}
\externaldocument{../08/08.tex}
\externaldocument{../09/09.tex}
\externaldocument{../10/10.tex}
\externaldocument{../11/11.tex}
\externaldocument{../12/12.tex}
\externaldocument{../13/13.tex}
\externaldocument{../14/14.tex}
%\externaldocument{../15/15.tex}
\externaldocument{../16/16.tex}
\externaldocument{../17/17.tex}
\externaldocument{../18/18.tex}
\externaldocument{../19/19.tex}

\setboolean{original}{false}

\begin{document}

\begin{center}
XV. Morphismes acycliques\\
M. Artin 
\end{center}

\selectlanguage{english}version : \showgitstatus\selectlanguage{french}
\tableofcontents

\else
\chapter{Morphismes acycliques}
%\addtocontents{toc}{M. Artin}
\begin{center}
M. Artin 
\end{center}
\fi

\setcounter{pageoriginal}{167}

% À faire :
% uniformiser notation pour la restriction F|X -> F_{|X} p. ex.
% rajouter des ~ pour éviter les termes mathématiques courts seuls en début de ligne.
% remplacer les \ref{...} par des (\ref{}) ; cf. p. ex. énoncé du 1.15
% ne pas mettre les remarques en italiques (cf. rem. 1.18).

\section*{Introduction}\pageoriginale

Soit $g:X\rightarrow Y$ un morphisme de schémas. Dans le premier
numéro, nous étudions des conditions pour que $g$ soit acyclique,
c'est-à-dire, tel que pour tout faisceau de torsion $F$ sur $Y$ on ait
$H^{q}(Y,F)\simeq H^{q}(X,g^{\ast}F)$, et que cela reste vrai si l'on
remplace $Y$ par un $Y'$ étale sur $Y$. Une condition naturelle
nécessaire d'acyclicité est que les fibres géométriques soient
acycliques, c'est-à-dire, aient une cohomologie triviale pour des
faisceaux de torsion constants. On verra que cette condition, jointe à
une condition locale, appelée acyclicité locale de $f$ (\hyperref[XV.def1.11]{1.11}), implique
que $f$ est acyclique.

Le théorème fondamental est qu'un morphisme lisse est localement
acyclique pour les faisceaux de torsion premiers aux caractéristiques
résiduelles \hyperref[XV.thm2.1]{2.1}, ce qui impliquera le théorème de
changement de base 
par un morphisme lisse \hyperrefext[XVI][XVI.thm1.1]{1.1}, qui sera développé, avec ses premières
conséquences, dans l'exposé suivant.

Le présent exposé est indépendant des exposés XI à XIV, et notamment
du « théorème de changement de base pour un morphisme propre ». Ce
dernier interviendra à nouveau de façon essentielle, mais en
conjonction avec le théorème fondamental du présent exposé, à partir
de l'exposé suivant.

\section{Généralités sur les morphismes globalement et localement
  acycliques}\etiquette[1]{XV.sec1}\pageoriginale

% g' n'est pas défini comme étant g x_Y Y' alors qu'il l'est en 1.3

\begin{enonce*}{Proposition 1.1}\etiquette[1.1]{XV.prop1.1}
Soit $g:X\rightarrow Y$ un morphisme de schémas. Les assertions
{\rm(i)} à {\rm(iii)} suivantes sont équivalentes:
\begin{itemize}
\item[(i)] Pour chaque $Y'\rightarrow Y$ étale (qu'on peut prendre affine
  au-dessus d'un ouvert affine de $Y$) et chaque faisceau d'ensembles
  $F$ sur $Y'$, le morphisme canonique $F\rightarrow
  g'_{\ast}{g'}^{\ast}F$ est injectif (\resp bijectif).

\item[(ii)] Pour chaque $Y'\rightarrow Y$ étale (qu'on peut prendre affine
  au-dessus d'un ouvert affine de $Y$) et chaque faisceau d'ensembles
  $F$ sur $Y'$, le morphisme canonique $H^{0}(Y',F)\rightarrow
  H^{0}(X',{g'}^{\ast}F)$ est injectif (\resp bijectif).

\item[(iii)] Pour chaque $Y'\rightarrow Y$ localement quasi-fini et
  chaque faisceau d'ensembles $F$ sur $Y'$, le morphisme canonique
  $H^{0}(Y',F)\rightarrow H^{0}(X',{g'}^{\ast}F)$ est injectif
  (\resp bijectif).

\item[(ii bis)] (Si $g$ est quasi-compact et quasi-séparé.) Comme
  {\rm(ii)}, mais en prenant $F$ faisceau constant de la forme
  $I_{Y'}$, où $I$ est un ensemble donné au préalable, avec $\card
  I\geq 2$.
\end{itemize}
\end{enonce*}

\noindent
\emph{Démonstration}. L'équivalence de (i) et (ii) est immédiate à
partir des définitions, et (iii) $\Rightarrow$ (ii) $\Rightarrow$ (ii
bis) sont triviales. On a (ii bis) $\Rightarrow$ (ii), car on peut
supposer $Y$ affine, donc $X$ et $Y$ quasi-compacts et quasi-séparés,
et on peut alors appliquer \hyperrefext[XII][XII.prop6.5]{6.5} (i). Pour l'implication (i)
$\Rightarrow$ (iii), notons d'abord le lemme suivant:

\begin{enonce*}{Lemme 1.2}\etiquette[1.2]{XV.lem1.2}
Soit\pageoriginale $f:Y'\rightarrow Y$ un morphisme localement
quasi-fini. Alors $f$ est « fini localement sur $Y'$ pour la
topologie étale », c'est-à-dire, on a des diagrammes commutatifs
$$\vcenter{
\xymatrix@=1.75cm{
Y'_i        \ar[r]^{a_i}   \ar[d]^{b_i}& Y'\ar[d]^{f}\\
Y_i          \ar[r]^{c_i}              & Y
}}\leqno{(*)}
$$
où les $a_i$, $c_i$ sont étales, les $b_i$ sont finis, et les $a_i$
forment un recouvrement de $Y'$ \footnote[1]{\Cf \EGA IV 18.12.1.}.
\end{enonce*}

En effet, soit $y'$ un point de $Y'$ et $y=f(y')$. Il suffit de
trouver un diagramme $({\ast})$ tel que $y'$ soit dans l'image de
$Y'_i$. Soit $\widetilde{Y}$ le spectre de l'anneau hensélisé de
$\Oo_{Y,y}$ \hyperrefext[VIII][VIII.subsec4.1]{4.1}
et $Z$ le localisé (ordinaire) de
$\widetilde{Y}'=\fprod{Y'}{\widetilde{Y}}{Y}$ en l'unique point
au-dessus de $y'$ et du point fermé $\widetilde{y}$ de
$\widetilde{Y}$. D'après \hyperrefext[VIII][VIII.subsec4.1]{4.1} (ii), % étiquette pour (ii) ?
$Z\rightarrow \widetilde{Y}$
est fini, et $Z$ est un ouvert induit de $\widetilde{Y}'$. Puisque
$\widetilde{Y}$ est limite de préschémas $Y_i$ étales au-dessus de
$Y$, on voit immédiatement \EGA IV 9 que pour $Y_i$ convenable il
existe un ouvert $Z_i$ de $\fprod{Y'}{Y_i}{Y}$, tel que
$Z=\fprod{Z_i}{\widetilde{Y}}{Y_i}$, et dans $({\ast})$ il suffira de
prendre $Y'_i=Z_i$, \cqfd


Supposons maintenant que (i) de \hyperref[XV.prop1.1]{1.1} soit vrai. Il est évident que
l'assertion (iii) de \hyperref[XV.prop1.1]{1.1} équivaut à l'assertion déduite de (i) en
remplaçant le morphisme $Y'\rightarrow Y$ étale par un morphisme
localement quasi-fini. Sous cette forme, l'assertion (pour $Y$, $Y'$
donnés) est locale sur $Y$ et $Y'$ pour la topologie étale, et on est
donc réduit par le lemme au cas d'un morphisme $f:Y'\rightarrow Y$
\emph{fini}.\pageoriginale Soit $f':X'\rightarrow X$ le morphisme
déduit de $f$ par le changement de base $g$. Alors le morphisme de
changement de base $g^{\ast}f_{\ast}F\rightarrow
{f'}_{\ast}{g'}^{\ast}F$ est bijectif \hyperrefext[VIII][VIII.cor5.6]{5.6}
D'après (i),
\sisi{le morphisme d'adjonction $f_{\ast}F\rightarrow g_{\ast}g^{\ast}f_{\ast}F$
est injectif (resp.~bijectif)}{le morphisme d'adjonction $f_{\ast}F\rightarrow g_{\ast}g^{\ast}f_{\ast}F$
est injectif (resp.~bijectif)}.
Par suite le morphisme $H^{0}(Y',F)\xrightarrow{\epsilon}
H^{0}(Y',g'_{\ast}{g'}^{\ast}F)$ de (iii) est le composé
$H^{0}(Y',F)=H^{0}(Y,f_{\ast}F)\xrightarrow{\epsilon'}H^{0}(Y,g_{\ast}g^{\ast}f_{\ast}F)=H^{0}(X,g^{\ast}f_{\ast}F)\simeq
H^{0}(X,{f'}_{\ast}{g'}^{\ast}F)=H^{0}(X',{g'}^{\ast}F)$.
Puisque $\epsilon'$ est injectif (\resp bijectif) il en est de même de
$\epsilon$. Cela achève la démonstration de \hyperref[XV.prop1.1]{1.1}.

\begin{enonce*}{Définition 1.3}\etiquette[1.3]{XV.def1.3}
On appelle morphisme \emph{$(-1)$-acyclique} (\resp
\emph{$0$-acyclique}) un morphisme $g:X\rightarrow Y$ ayant une des
propriétés équivalentes {\rm (i)} à {\rm (iii)} (\resp {\rm(i) à
  (iii)} respées) de $1.1$. Le morphisme est dit
\emph{universellement} $(-1)$-acyclique (\resp $0$-acyclique) si pour
chaque morphisme de changement de base $Y'\rightarrow Y$, le morphisme
$g'=\fprod{g}{Y'}{Y}:\fprod{X}{Y'}{Y}\rightarrow Y'$ est
$(-1)$-acyclique (\resp $0$-acyclique). Enfin, un schéma $X$ est dit
$(-1)$-acyclique (\resp $0$-acyclique) s'il est non-vide (\resp
connexe et non vide).
\end{enonce*}

\begin{enonce*}{Corollaire 1.4}\etiquette[1.4]{XV.cor1.4}
Soit $g:X\rightarrow Y$ un morphisme $0$-acyclique. Alors pour chaque
$Y'\rightarrow Y$ localement quasi-fini et chaque faisceau en groupes
$F$ sur $Y'$, le morphisme canonique
$$
H^{1}(Y',F)\longrightarrow H^{1}(X',{g'}^{\ast}F)
$$
est injectif.
\end{enonce*}

Cela résulte aussitôt de \hyperrefext[XII][XII.prop6.5]{6.5} (i).

\begin{enonce*}{Corollaire 1.5}\etiquette[1.5]{XV.cor1.5}
Soit\pageoriginale $g:X\rightarrow Y$ un morphisme. Si $g$ est
surjectif (\ie ses fibres sont non vides, \ie ses fibres sont
$(-1)$-acycliques) alors $g$ est $(-1)$-acyclique. La réciproque est
vraie si $g$ est quasi-compact et quasi-séparé.
\end{enonce*}

La première assertion est évidente sur la forme \hyperref[XV.prop1.1]{1.1} (ii) par
considération des fibres des faisceaux envisagés. La deuxième résulte
de la définition sous la forme \hyperref[XV.prop1.1]{1.1} (iii) et du sorite de
passage à la limite \hyperrefext[VII][VII.sec5]{5},
compte tenu que pour tout $y \in Y$, $y=\Spec k(y)$ est
limite projective de schémas affines localement quasi-finis sur $Y$,
savoir les voisinages affines de $y$ dans $\bar{\{y\}}$ (muni de la
structure réduite induite). \sisi{}{--} Notons que \hyperref[XV.cor1.5]{1.5} implique que lorsque $g$
est quasi-compact et quasi-séparé, alors $g$ $(-1)$-acyclique implique
déjà $g$ \emph{universellement} $(-1)$-acyclique.

\begin{enonce*}{Proposition 1.6}\etiquette[1.6]{XV.prop1.6}
Soient $g:X\rightarrow Y$ un morphisme de schémas, $\LL\subset
\PP$ un ensemble de nombres premiers, et $n \in \NN$ un
entier naturel (\resp $n=1$). Les assertions {\rm (i)} à {\rm(iii)}
(\resp les assertions {\rm(i)} à {\rm(iii)} respées) suivantes sont
équivalentes. Si de plus $g$ est quasi-compact et quasi-séparé, alors
{\rm(i)} à {\rm(iii)} (\resp {\rm(i)} à {\rm(iii)} respées) sont aussi
équivalentes à {\rm(iv)} (\resp {\rm(iv)} respée).
\begin{enumerate}
\renewcommand{\theenumi}{\roman{enumi}}
\renewcommand{\labelenumi}{(\theenumi)}
\item Pour chaque $Y'\rightarrow Y$ étale et chaque faisceau abélien
  de $\LL$-torsion (\resp de $\text{\rm ind}$-$\LL$-groupes)
  \emph{\hyperrefext[IX][IX.def1.5]{1.5}} $F$ sur $Y'$, on a
$$
\begin{cases}
F\xrightarrow{\sim} {g'}_{\ast}{g'}^{\ast}F, &\\
(R^{q}{g'}_{\ast}){g'}^{\ast}F=0 & \mbox{si}\quad 1\leq q\leq n\quad
(\resp \mbox{si\; } q=1){\quoi}.
\end{cases}
$$

\item Pour\pageoriginale chaque $Y'\rightarrow Y$ étale et chaque
faisceau abélien de $\LL$-torsion (\resp de
ind-$\LL$-groupes) $F$ sur $Y'$, le morphisme
canonique
$$
H^{q}(Y',F)\longrightarrow H^{q}(X',{g'}^{\ast}F)
$$
est bijectif si $q\leq n$ et injectif si $q=n+1$ (\resp est bijectif
si $q\leq 1$).

\item Pour chaque $Y'\rightarrow Y$ localement quasi-fini et chaque
faisceau abélien de $\LL$-torsion (\resp de
ind-$\LL$-groupes) $F$ sur $Y'$, le morphisme
$$
H^{q}(Y',F)\longrightarrow H^{q}(X', {g'}^{\ast}F)
$$
est bijectif si $q\leq n$ et injectif si $q=n+1$ (\resp est bijectif
si $q\leq 1$).

\item Le morphisme $g$ est surjectif, et pour chaque $Y'\rightarrow Y$
localement quasi-fini, chaque $\nu>0$, et chaque $\ell\in \LL$,
le morphisme canonique
$$
H^{q}(Y',\ZZ/\ell^{\nu})\longrightarrow H^{q}(X', \ZZ/\ell^{\nu})
$$
est \emph{surjectif} pour $0\leq q\leq n$ (\resp pour chaque
$Y'\rightarrow Y$ localement quasi-fini et chaque $\LL$-groupe
fini ordinaire $G$, le morphisme canonique
$$
H^{i}(Y',G)\longrightarrow H^{i}(X', G)
$$
est \emph{surjectif} pour $i=0,1$). 
\end{enumerate}
\end{enonce*}

\noindent
\emph{Démonstration de}\pageoriginale \hyperref[XV.prop1.6]{1.6}. Les implications (ii)
$\Leftarrow$ (iii) 
$\Rightarrow$ (iv) sont triviales.
% surjectivité triviale ? (g n'est pas nécessairement cohérent)
Or $(R^{q}{g'}_{\ast}){g'}^{\ast}F$
est le faisceau associé au préfaisceau
$R^{q}(Y'')=H^{q}(X'',{g''}^{\ast}F)$ ($Y''\rightarrow Y'$ étale). Si
(ii) est vrai, on a $H^{q}(Y'',F)\xrightarrow{\sim}
H^{q}(X'',{g''}^{\ast}F)$, donc le faisceau associé est nul\sisi{
si $1 \leq q \leq n$}{}, \ie (ii) $\Rightarrow$ (i). L'implication (i) $\Rightarrow$ (ii) résulte de 1.1
(i) (qui correspond au cas $n=0$) joint à la suite spectrale de Leray
$$
H^{p}(Y',(R^{q}{g'}_{\ast}){g'}^{\ast}F)\Longrightarrow
H^{p+q}(X',{g'}^{\ast}F)
$$
dans le cas abélien, et à la suite exacte
$$
0\longrightarrow H^{1}(Y',{g'}_{\ast}{g'}^{\ast}F)\longrightarrow
H^{1}(X',{g'}^{\ast}F)\longrightarrow
H^{0}(Y',(R^{1}{g'}_{\ast}){g'}^{\ast}F){\quoi}.
$$
dans le cas des faisceaux de groupes.

(i) $\Rightarrow$ (iii): Puisque (i) et (ii) sont équivalentes, il est
clair que (iii) équivaut à l'énoncé qui est déduit de (i) en remplaçant
le mot « étale »  par les mots « localement quasi-fini ». Sous cette forme c'est, pour $Y$, $Y'$ donnés, une assertion
locale sur $Y'$ et sur $Y$ pour la topologie étale, et on est donc
réduit par \hyperref[XV.lem1.2]{1.2} au cas où $f:Y'\rightarrow Y$ est un morphisme
\emph{fini}. Soit $f':X'\rightarrow X$ le morphisme
\sisi{déduit de~$f$ par le changement de base~$g$}{déduit de~$f$ par le changement de base~$g$}.
On a d'après \hyperrefext[VIII][VIII.prop5.5]{5.5} (\resp \hyperrefext[VIII][VIII.rem5.8]{5.8}) (appliqué trois fois)
$$
f_{\ast}(R^{q}{g'}_{\ast}){g'}^{\ast}F\simeq
R^{q}(fg')_{\ast}{g'}^{\ast}F\simeq
R^{q}(gf')_{\ast}{g'}^{\ast}F\simeq
(R^{q}g_{\ast}){f'}_{\ast}{g'}^{\ast}F\simeq
(R^{q}g_{\ast})g^{\ast}(f_{\ast}F)
$$
et ce dernier est nul si (i) est vrai pour les valeurs de $q$
envisagées, d'où (i) $\Rightarrow$ (iii).

Il reste à démontrer (iv) $\Rightarrow$ (iii); or ceci résulte
aussitôt de\pageoriginale \hyperrefext[XII][XII.prop6.5]{6.5}.

\begin{enonce*}{Définition 1.7}\etiquette[1.7]{XV.def1.7}
Soient $n\in \NN$ et $\LL\subset \PP$. Un
morphisme $g:X\rightarrow Y$ qui satisfait à une des conditions
équivalentes {\rm(i)} à {\rm(iii)} (\resp {\rm(i)} à {\rm(iii)}
respées) de \hyperref[XV.prop1.6]{1.6} est appelé \emph{morphisme $n$-acyclique}
(\resp \emph{morphisme $1$-asphérique}) pour $\LL$. On dit que $g$ est \emph{acyclique}
pour $\LL$ s'il est $n$-acyclique pour $\LL$ pour chaque
$n$. On dit que $g$ est \emph{universellement $n$-acyclique} (\resp
\emph{universellement $1$-asphérique}) pour $\LL$, si pour
chaque $Y'\rightarrow Y$ le morphisme $g'=\fprod{g}{Y'}{Y}$ est
$n$-acyclique (\resp $1$-asphérique) pour $\LL$. Enfin, un
préschéma $X$ est dit $n$-\emph{acyclique} (\resp
$1$-\emph{asphérique}) pour $\LL$, s'il est $0$-acyclique $1.3$
et si pour chaque $\ell \in \LL$ et chaque $1\leq q\leq n$ on a
$H^{q}(X,\ZZ/\ell)=0$ (\resp et pour chaque $\LL$-groupe
fini ordinaire $G$, on a $H^{1}(X,G)=0$).
\end{enonce*}

\begin{enonce*}[remark]{Remarques 1.8}\etiquette[1.8]{XV.rem1.8}
\begin{enumerate}
\renewcommand{\theenumi}{\alph{enumi}}
\renewcommand{\labelenumi}{\theenumi)}
\item Si $g$ est quasi-compact et quasi-séparé, il suffit, pour que
$g$ soit universellement $n$-acyclique, que $g'$ soit $n$-acyclique
chaque fois que $Y'\rightarrow Y$ est localement de présentation
finie. Nous n'avons pas besoin de ce fait et en laissons la
démonstration, qui est un exercice de passage à la limite, au lecteur.

\item Nous ignorons si un morphisme $n$-acyclique pour $\LL$
est déjà universellement $n$-acyclique pour $\LL$. La même
question se pose pour les variantes locales plus bas \hyperref[XV.def1.11]{1.11}.

\item Lorsque $n=0$, on retrouve la notion de \hyperref[XV.def1.3]{1.3}. Bien
entendu, on aurait pu rédiger \hyperref[XV.prop1.6]{1.6}. et \hyperref[XV.def1.7]{1.7}. de
façon à inclure le cas $n=-1$. La proposition qui suit est
également valable dans ce cas.\pageoriginale
\end{enumerate}
\end{enonce*}

\begin{enonce*}{Proposition 1.9}\etiquette[1.9]{XV.prop1.9}
\begin{enumerate}
\renewcommand{\theenumi}{\roman{enumi}}
\renewcommand{\labelenumi}{(\theenumi)}
\item Soient $X\xrightarrow{g}Y\xrightarrow{h}Z$ des morphismes de
schémas, $\LL\subset \PP$, et $n\in
\NN$. Supposons que $g$ soit $0$-acyclique (\resp
$1$-asphérique pour $\LL$, \resp $n$-acyclique pour
$\LL$). Soit $f=hg$. Alors pour tout faisceau d'ensembles
(\resp de ind\sisi{-}{ }$\LL$-groupes, \resp abélien de
$\LL$-torsion) $F$ sur $Z$, on a un isomorphisme
$f_{\ast}f^{\ast}F\xrightarrow{\sim} h_{\ast}h^{\ast}F$ (\resp des
isomorphismes $(R^{p}f_{\ast})f^{\ast}F\xrightarrow{\sim}
(R^{p}h_{\ast})h^{\ast}F$ si $p\leq 1$, \resp si $p\leq n$). En
particulier, $h$ est $0$-acyclique (\resp $1$-asphérique pour
$\LL$, \resp $n$-acyclique pour $\LL$) si et seulement
si $f$ l'est.

\item Soit $\{g_{\alpha}:X_{\alpha}\rightarrow Y_{\alpha}\}$ un
système projectif de morphismes de schémas tels que les morphismes
$X_{\alpha}\rightarrow X_{\beta}$ et $Y_{\alpha}\rightarrow Y_{\beta}$
soient affines et que les $g_{\alpha}$ soient quasi-compacts et
quasi-séparés. Soit $g:X\rightarrow Y$ la limite projective des $g_{\alpha}$
\hyperrefext[VII][VII.subsec5.1]{5.1}.
Alors si tous les $g_{\alpha}$ sont $0$-acycliques (\resp
$1$-asphériques pour $\LL$, \resp $n$-acycliques pour
$\LL$), il en est de même de $g$.
\end{enumerate}
\end{enonce*}

\noindent
\emph{Démonstration}. Pour (i), on a d'abord
$$
h_{\ast}h^{\ast}F\xrightarrow{\sim}
h_{\ast}(g_{\ast}g^{\ast})h^{\ast}F\simeq f_{\ast}f^{\ast}F{\quoi}.
$$
La deuxième assertion de (i) est conséquence immédiate de la suite
spectrale de Leray
$$
(R^{p}h_{\ast})(R^{q}g_{\ast})f^{\ast}F\Longrightarrow
(R^{p+q}f_{\ast})f^{\ast}F{\quoi}.
$$
Enfin l'assertion non-abélienne résulte de la suite exacte \hyperrefext[XII][XII.prop3.2]{3.2}.
L'assertion\pageoriginale (ii) résulte facilement de la théorie
de passage à la limite \hyperrefext[VII][VII.sec5]{5} et VI. % étiquette pour chapitre ?

\begin{enonce*}{Corollaire 1.9.1}\etiquette[1.9.1]{XV.cor1.9.1}
Soit $g:X\rightarrow Y$ un morphisme quasi-compact et quasi-séparé, et
$n\in \NN$, $\LL\subset \PP$. Si $g$ est
$n$-acyclique pour $\LL$ (\resp $1$-asphérique pour
$\LL$), alors pour tout point géométrique $\bar{y}$ de $Y$,
algébrique sur un point $y$ de $Y$, la fibre géométrique $X_{\bar{y}}$
est $n$-acyclique pour $\LL$ (\resp $1$-asphérique pour $\LL$).
\end{enonce*}

En effet, $\bar{y}$ est limite projective de schémas affines $Y_i$ qui
sont localement quasi-finis sur $Y$, et comme les
$X_i=\fprod{X}{Y_i}{Y}\rightarrow Y_i$ sont $n$-acycliques pour
$\LL$ (\resp\ldots), il en est de même de
$X_{\bar{y}}\rightarrow \bar{y}$ en vertu de \hyperref[XV.prop1.9]{1.9} (ii), d'où la
conclusion annoncée.

\begin{enonce*}{Proposition 1.10}\etiquette[1.10]{XV.prop1.10}
Soient $g:X\rightarrow Y$ un morphisme de schémas, $\LL\subset
\PP$, et $n\in \NN$. Alors les conditions {\rm(i)} à
{\rm(iv)} ci-dessous sont équivalentes.
\begin{enumerate}
\renewcommand{\theenumi}{\roman{enumi}}
\renewcommand{\labelenumi}{(\theenumi)}
\item Soit $\bar{y}$ un point géométrique de $Y$ et $\bar{x}$ un point
géométrique de $X$ au-dessus de $\bar{y}$. Soient $\widetilde{X}$,
$\widetilde{Y}$ les localisés stricts de $X$, $Y$ aux points
$\bar{x}$, $\bar{y}$, et soit $\widetilde{g}:\widetilde{X}\rightarrow
\widetilde{Y}$ le morphisme induit par $g$. Alors $\widetilde{g}$ est
$0$-acyclique (\resp $1$-asphérique pour $\LL$, \resp
$n$-acyclique pour $\LL$).

\item Pour chaque diagramme à carrés cartésiens
$$
\vcenter{
\xymatrix@C=2cm@R=1.25cm{
X   \ar[d]^{g}& X'\ar[l]_{j}  \ar[d]^{g'} & X'' \ar[l]_{j'} \ar[d]^{g''}\\
Y    & Y'\ar[l]_{i} &   Y'' \ar[l]_{i'} 
}}\leqno{(\ast)}
$$
avec\pageoriginale $i$ étale et $i'$ étale de présentation finie, et
chaque faisceau d'ensembles (\resp de ind-$\LL$-groupes, \resp
abélien de $\LL$-torsion) $F$ sur $Y''$, le morphisme de
changement de base \hyperrefext[XII][XII.eq4.1]{4.1}, \hyperrefext[XII][XII.eq4.2]{4.2}
$$
{g'}^{\ast}i'_{\ast}F\longrightarrow {j'}_{\ast}{g''}^{\ast}F
$$
est bijectif (\resp les morphismes
$$
{g'}^{\ast}(R^{q}i'_{\ast})F\longrightarrow
(R^{q}j'_{\ast}){g''}^{\ast}F
$$
sont bijectifs pour $q=0$, $1$, \resp sont bijectifs pour $0\leq q\leq
n$ et injectifs pour $q=n+1$).

\item Même énoncé que {\rm(ii)}, sauf que le morphisme $i'$ est
supposé seulement quasi-fini et quasi-séparé.

\item Même énoncé que {\rm(ii)}, mais en supposant $i$ seulement
localement quasi-fini, et en revanche $i'$ une immersion ouverte
quasi-compacte.
\end{enumerate}
\end{enonce*}

\noindent
\emph{Démonstration}. L'implication (iii) $\Rightarrow$ (ii) est
triviale. Vérifions que (i) $\Rightarrow$ (iii). Il suffit de regarder
le morphisme de (iii) fibre par fibre. Soient $\bar{y}$ un point
géométrique de $Y'$ et $\bar{x}$ un point géométrique de $X'$
au-dessus de $\bar{y}$. Soit
$$
\xymatrix@=1.75cm{
\widetilde{X}'   \ar[d]^{\widetilde{g}'}& \widetilde{X}''
\ar[l]_{\widetilde{j}'} \ar[d]^{\widetilde{g}''}\\
 \widetilde{Y}'           & \widetilde{Y}'' \ar[l]_{\widetilde{i}'}
}
$$
le diagramme cartésien déduit de $({\ast})$, où $\widetilde{X}'$,
$\widetilde{Y}'$ sont les localisés stricts des $X'$, $Y'$ aux points
$\bar{x}$, $\bar{y}$ et où
$\widetilde{Y}''=\fprod{Y''}{\widetilde{Y}'}{Y'}$,
$\widetilde{X}''=\fprod{X''}{\widetilde{X}'}{X'}$.\pageoriginale On a
\hyperrefext[VIII][VIII.thm5.2]{5.2} et \hyperrefext[VIII][VIII.rem5.3]{5.3}
$$
H^{q}(\widetilde{Y}'',\widetilde{F})\simeq
(R^{q}i'{}_{\ast}F)_{\bar{y}}\simeq
({g'}^{\ast}(R^{q}i'_{\ast})F)_{\bar{x}}{\quoi},
$$
où $\widetilde{F}$ est l'image inverse de $F$ sur $\widetilde{Y}''$,
et
$$
H^{q}(\widetilde{X}'',{\widetilde{g}''}{}^{\ast}\widetilde{F})\simeq
((R^{q}j'_{\ast}){g''}{}^{\ast}F)_{\bar{x}}{\quoi}.
$$
Puisque $\widetilde{i}'$ est quasi-fini, on peut appliquer \hyperref[XV.prop1.6]{1.6} (iii)
(\resp \hyperref[XV.cor1.5]{1.5} (iii)) au morphisme
$$
H^{q}(\widetilde{Y}'',\widetilde{F})\longrightarrow
H^{q}(\widetilde{X}'',{\widetilde{g}''}{}^{\ast}\widetilde{F})
$$
pour les valeurs de $q$ envisagées, d'où (i) $\Rightarrow$ (iii).

(ii) $\Rightarrow$ (i). Supposons (ii) vrai et vérifions le critère
\hyperref[XV.prop1.6]{1.6} (ii) d'acyclicité pour un morphisme
$\widetilde{g}:\widetilde{X}\rightarrow \widetilde{Y}$. Soit
$\widetilde{Y}'\rightarrow \widetilde{Y}$ étale, avec $\widetilde{Y}'$
affine, et soit $\widetilde{F}$ un faisceau d'ensembles (\resp abélien
de $\LL$-torsion, \resp de ind-$\LL$-groupes) sur
$\widetilde{Y}'$. On doit démontrer que
$$
H^{q}(\widetilde{Y}',\widetilde{F})\xrightarrow{\sim}
H^{q}(\widetilde{X}', \widetilde{g}'{}^{\ast}\widetilde{F}){\quoi},
$$
où $\widetilde{g}':\widetilde{X}'\rightarrow \widetilde{Y}'$ est le
morphisme déduit de $\widetilde{g}$ par changement de base. Or
écrivons $\widetilde{Y}=\varprojlim Y_{\alpha}$ où les $Y_{\alpha}$
sont étales et affines au-dessus de $Y$ (supposant $Y$ affine, ce qui
est loisible). On peut « descendre »  le morphisme étale
$\widetilde{Y}'\rightarrow \widetilde{Y}$ à un des $Y_{\alpha}$,
disons en $Y'_{0}\rightarrow Y_{0}$, et on a
$\widetilde{Y}'=\varprojlim
Y'_{\alpha}=\fprod{Y}{Y'_{0}}{Y_{0}}$. Or chaque faisceau
$\widetilde{F}$ sur $\widetilde{Y}'$ est limite inductive des
faisceaux ${f'_{\alpha}}^{\ast}{f'_{\alpha}}_{\ast}\widetilde{F}$, où
$f'_{\alpha}:\widetilde{Y}'_{\alpha}\rightarrow Y'_{\alpha}$ est le
morphisme canonique. D'après \sisi{XII 1.2 (v) et 1.6 (iii)}{\hyperrefext[IX][IX.prop1.2]{1.2} et \hyperrefext[IX][IX.prop1.6]{1.6}} les faisceaux
$f'_{\alpha}{}_{\ast}F$ sont des faisceaux d'ensembles (\resp abéliens
de $\LL$-torsion, \resp de
ind-$\LL$-groupes). Compte\pageoriginale tenu de \hyperrefext[VII][VII.prop3.3]{3.3},
on voit donc qu'il suffit de traiter le cas où $\widetilde{F}$ se descend
à un des $Y'_{\alpha}$, disons à un $F_{0}$ sur $Y'_{0}$. En
appliquant (ii) au diagramme
$$
\xymatrix@C=2.25cm@R=1.5cm{
X \ar[d]_{g}& X_0 \ar[l] \ar[d]_{g_0}& X'_0 \ar[l]_{j'_0} \ar[d]_{g'_0}\\
Y & Y_0 \ar[l]& Y'_0 \ar[l]_{i'_0},
}
$$
on déduit que
${g_{0}}^{\ast}(R^{q}{i'_0}_{\ast})F_{0}\xrightarrow{\sim}
(R^{q}{j'_0}_{\ast}){g'_0}^{\ast}F_{0}$ pour les
valeurs de $q$ applicables. Or les points géométriques $\bar{y}$,
$\bar{x}$ se relèvent canoniquement en des points géométriques de
$Y_{0}$, $X_{0}$ (nous les notons par les mêmes symboles), et
on a \hyperrefext[VIII][VIII.thm5.2]{5.2} et \hyperrefext[VIII][VIII.rem5.3]{5.3}
$$
H^{q}(\widetilde{Y}',\widetilde{F})=((R^{q}{i'_0}_{\ast})F_{0})_{\bar{y}}=({g_0}^{\ast}
(R^{q}{i'_0}_{\ast})F_0)_{\bar{x}}
$$
et
$$
H^{q}(\widetilde{X}',{\widetilde{g}'}{}^{\ast}\widetilde{F})=((R^{q}j'_{0\ast}){g'_0}^{\ast}F_{0})_{\bar{x}}
$$
pour les valeurs de $q$ envisagées, d'où le résultat.

Reste à montrer que les conditions (i) à (iii) sont équivalentes à
(iv). Or, utilisant la forme (i), on voit que ces conditions sont
stables par extension de la base $Y'\rightarrow Y$ localement
quasi-finie, d'où résulte aussitôt, sous la forme (ii), qu'elles
impliquent (iv). D'autre part, (iv) implique (iii), car pour vérifier
(iii) on voit tout de suite, utilisant par exemple la suite spectrale
de Leray, ou le résultat de descente \hyperrefext[XII][XII.prop6.8]{6.8}, que l'on peut supposer
$Y'$ et $Y''$ affines. Mais alors par le « Main Theorem »  \EGA IV
8.12.8 $Y''\rightarrow Y'$ se factorise en
$Y''\xrightarrow{i'_1}Y'_1\xrightarrow{i'_2}Y'$, avec~$i'_1$
une immersion ouverte,\pageoriginale et~$i'_2$ un morphisme
\emph{fini}. Utilisant \hyperrefext[VIII][VIII.prop5.5]{5.5}, \hyperrefext[VIII][VIII.rem5.8]{5.8} pour~$i'_2$, on est réduit à
vérifier (ii) pour~$i'_1$, qui relève de (iv). Ceci achève la
démonstration de~\hyperref[XV.prop1.10]{1.10}.

\begin{enonce*}{Définition 1.11}\etiquette[1.11]{XV.def1.11}
Un morphisme $g:X\rightarrow Y$ de préschémas qui satisfait à une des
conditions équivalentes {\rm(i)} à {\rm(iv)} de $1.10$ est appelé
morphisme \emph{localement $0$-acyclique} (\resp \emph{localement
$1$-asphérique} pour $\LL$, \resp \emph{localement
$n$-acyclique} pour $\LL$). On définit d'une manière évidente
(comparer \hyperref[XV.def1.7]{1.7}) les notions de morphisme \emph{localement acyclique},
et de morphisme \emph{universellement localement $0$-acyclique}
(\resp\ldots) pour~$\LL$.
\end{enonce*}

On peut naturellement varier les énoncés (i) à (iii) de \hyperref[XV.prop1.10]{1.10}. Ainsi,
dans (iii) et (iv) on peut se borner à un faisceau $F$ constant et de
la forme $G_{I}$,
% quésako cette notation ?
 $G$ un ensemble fini (\resp un $\LL$-groupe
fini, \resp $G$ de la forme $\ZZ/\ell^{\nu}\ZZ$, avec
$\ell\in\LL$, $\nu\in \NN$); cela résulte facilement de
\hyperrefext[VIII][VIII.thm5.2]{5.2}, \hyperrefext[VIII][VIII.rem5.3]{5.3} et de \hyperrefext[XII][XII.prop6.5]{6.5}.
Signalons également la variante suivante:

\begin{enonce*}{Corollaire 1.12}\etiquette[1.12]{XV.cor1.12}
Supposons $g:X\rightarrow Y$ localement de type fini. Pour vérifier la
$0$-acyclicité (\resp\ldots) locale, il suffit de vérifier $1.10$
{\rm(i)} dans le cas où le point géométrique $\bar{x}$ est
\emph{fermé} dans la fibre géométrique $X_{\bar{y}}$.
\end{enonce*}

En effet, dans la démonstration de (i) $\Rightarrow$ (iii) de \hyperref[XV.prop1.10]{1.10}, on
a vérifié (iii) fibre par fibre, et il suffisait de le faire pour les
fibres aux points géométriques $\bar{x}$ qui sont fermés dans leur
fibre géométrique,\pageoriginale puisque $g$ est localement de type
fini \hyperrefext[VIII][VIII.rem3.13]{3.13} b).

\begin{enonce*}{Proposition 1.13}\etiquette[1.13]{XV.prop1.13}
\begin{enumerate}
\renewcommand{\theenumi}{\roman{enumi}}
\renewcommand{\labelenumi}{(\theenumi)}
\item Soient $X\xrightarrow{g}Y\xrightarrow{h}Z$ des morphismes de
schémas, $\LL\subset \PP$, et $n\in \NN$, et
supposons que $g$ soit surjective, et localement $0$-acyclique (\resp
localement $1$-asphérique pour $\LL$, \resp localement
$n$-acyclique pour $\LL$). Alors $h$ est localement
$0$-acyclique (\resp\ldots) si et seulement si $f=hg$ l'est.

\item Soit $\{g_{\alpha}:X_{\alpha}\rightarrow Y_{\alpha}\}$ un
système projectif de morphismes de préschémas tels que les morphismes
$X_{\alpha}\rightarrow X_{\beta}$ et $Y_{\alpha}\rightarrow Y_{\beta}$
soient affines. Soit $g$ la limite projective des $\{g_{\alpha}\}$
\hyperrefext[VII][VII.subsec5.1]{5.1}.
Alors si tous les $g_{\alpha}$ sont localement
$0$-acycliques (\resp\ldots), il en est de même de $g$.
\end{enumerate}
\end{enonce*}

C'est un sorite analogue à \hyperref[XV.prop1.9]{1.9}.

\begin{enonce*}[remark]{Remarques 1.14}\etiquette[1.14]{XV.rems1.14}
Pour simplifier sa tâche, le rédacteur a omis dans \hyperref[XV.prop1.10]{1.10} 
à \hyperref[XV.prop1.13]{1.13} de
traiter également le cas de $n=-1$, et il laisse au lecteur le soin de
se convaincre que les énoncés, définitions et démonstrations
s'appliquent essentiellement sans changement à ce cas. Ainsi, avec les
notations de \hyperref[XV.prop1.10]{1.10}, la condition (i) s'énonce en disant que les
morphismes induits locaux $\widetilde{X}\rightarrow \widetilde{Y}$
sont $(-1)$-acycliques i.e. \hyperref[XV.cor1.5]{1.5} \emph{surjectifs}, et les conditions (ii)
à (iv) s'énoncent en disant que le morphisme de changement de base
${g'}^{\ast}({i'}_{\ast}(F))\rightarrow {j'}_{\ast}({g''}^{\ast}(F))$
est un \emph{monomorphisme}. Sous ces conditions, on dira donc que $g$
est \emph{localement $(-1)$-acyclique}. Cette condition est donc
stable\pageoriginale par extension localement quasi-finie de la
base. Comme exemple, signalons qu'un morphisme plat est localement
$(-1)$-acyclique, car pour un tel morphisme la condition de
surjectivité pour les $\widetilde{X}\rightarrow \widetilde{Y}$ est une
conséquence bien connue de la platitude de ces morphismes. D'autre
part, signalons aussi que lorsque $f$ est localement de présentation
finie, alors $f$ est localement $(-1)$-acyclique si et seulement si il
est \emph{universellement ouvert}, du moins lorsque $Y$ est localement
noethérien: cela résulte en effet aisément de \EGA IV 14.4.9,
1.10.3. Pour un tel morphisme, la $(-1)$-acyclicité locale implique
par suite la $(-1)$-acyclicité locale \emph{universelle}. (Il est
d'ailleurs très probable \footnote[1]{Cela a été effectivement vérifié
par M. Raynaud. \Cf réédition de \EGA IV 14!} que dans ces derniers
énoncés, l'hypothèse localement noethérienne sur $Y$ soit inutile:
c'est en tous cas ainsi si l'ensemble des composantes irréductibles de
$Y$ est localement fini \EGA IV 14.4.10.)
\end{enonce*}

\begin{enonce*}{Théorème 1.15}\etiquette[1.15]{XV.thm1.15}
Soient $g:X\rightarrow Y$ un morphisme quasi-compact et quasi-séparé,
$\LL\subset \PP$ un ensemble de nombres premiers, $n$ un
entier naturel (\resp $n=1$). On suppose que $g$ est $(n-1)$-acyclique
pour $\LL$ et localement $(n-1)$-acyclique pour $\LL$
$1.3$, $1.7$, $1.11$, $1.14$. Soit $F$ un faisceau sur $Y$, que pour
$n\geq 1$ on suppose muni d'une structure de faisceau abélien de
$\LL$-torsion (\resp d'une structure de faisceau en groupes),
et soit $\xi$ un élément de $H^{n}(X,g^{\ast}(F))$. Pour que $\xi$
provienne d'un élément de $H^{n}(Y,F)$ (qui est alors uniquement
déterminé, grâce à l'hypothèse de $(n-1)$-acyclicité sur $g$), il faut
et il suffit que pour tout point géométrique $\bar{y}$ de $Y$,
algébrique sur un point $y$ de $Y$, l'image $\xi_{\bar{y}}$ de $\xi$
dans $H^{n}(X_{\bar{y}},F|X_{\bar{y}})$ soit dans l'image de
$H^{n}(\bar{y},F_{\bar{y}})$ (ce qui, pour $n\geq 1$,
signifie\pageoriginale qu'on a $\xi_{\bar{y}}=0$).
\end{enonce*}

Notons tout de suite qu'on conclut de \hyperref[XV.thm1.15]{1.15} que, sous les conditions
préliminaires envisagées pour $g$, $g$ est $n$-acyclique pour
$\LL$ (\resp $1$-asphérique pour $\LL$) si et seulement
si pour tout $\bar{y}$ comme dessus, $X_{\bar{y}}$ est $n$-acyclique
pour $\LL$ (\resp est $1$-asphérique pour $\LL$): le « si »  résulte en effet directement de \hyperref[XV.thm1.15]{1.15} et des définitions, le
« seulement si »  étant de toutes façons immédiat par
\hyperref[XV.cor1.9.1]{1.9.1}. Ceci dit, une récurrence immédiate sur $n$ fournit le

\begin{enonce*}{Corollaire 1.16}\etiquette[1.16]{XV.cor1.16}
Soit $g$ un morphisme quasi-compact et quasi-séparé qui est
localement $(n-1)$-acyclique pour $\LL$, où $n$ est un entier
naturel donné (\resp où $n=1$). Alors $g$ est $n$-acyclique pour
$\LL$ (\resp est $1$-asphérique pour $\LL$) si et
seulement si pour tout point géométrique $\bar{y}$ de $Y$, algébrique
sur un point $y$ de $Y$, la fibre $X_{\bar{y}}$ est $n$-acyclique pour
$\LL$ (\resp $1$-asphérique pour $\LL$).
\end{enonce*}

\begin{enonce*}{Corollaire 1.17}\etiquette[1.17]{XV.cor1.17}
Soient $g:X\rightarrow Y$ un morphisme de préschémas,
$\LL\subset \PP$ un ensemble de nombres premiers, $n$ un
entier naturel (\resp $n=1$). Pour que $g$ soit localement
$n$-acyclique pour $\LL$ (\resp localement $1$-asphérique pour
$\LL$) il faut et il suffit que pour tout point géométrique
$\bar{x}$ de $X$, si $\widetilde{g}:\widetilde{X}\rightarrow
\widetilde{Y}$ est le morphisme des localisés stricts de $X$ et $Y$
(en $\bar{x}$ et le point géométrique correspondant $\bar{y}$ de $Y$)
induit par $g$, chaque fibre géométrique $X_{\bar{z}}$ de $g$,
relativement à un point géométrique $\bar{z}$ de $Y$ algébrique sur un
point $z$ de $Y$, soit $n$-acyclique pour $\LL$ (\resp
$1$-asphérique pour $\LL$).
\end{enonce*}

La nécessité résulte en effet du critère \hyperref[XV.prop1.10]{1.10}. 
(i) et de \hyperref[XV.cor1.9.1]{1.9.1}.,
pour\pageoriginale la suffisance on procède par récurrence sur $n$, ce
qui nous permet, lorsque $n\geq 1$, de supposer que $g$ est déjà
localement $(n-1)$-acyclique pour $\LL$. Il en est alors de
même des morphismes $\widetilde{g}$, et \hyperref[XV.cor1.16]{1.16} implique alors que
$\widetilde{g}$ est $n$-acyclique (\resp $1$-asphérique pour
$\LL$), ce qui en vertu du critère \hyperref[XV.prop1.10]{1.10} (i) implique que $g$
est localement $n$-acyclique (\resp localement $1$-asphérique) pour
$\LL$.

\begin{enonce*}{Corollaire 1.18}\etiquette[1.18]{XV.cor1.18}
Avec les notations de \hyperref[XV.cor1.17]{1.17}, si $g$ est localement de type fini,
alors dans le critère énoncé, on peut se borner à prendre des points
géométriques $\bar{x}$ dont le point $x$ est fermé dans sa fibre.
\end{enonce*}

Cela résulte en effet de la démonstration donnée et de \hyperref[XV.cor1.12]{1.12}.

\noindent
\emph{Démonstration de} \hyperref[XV.thm1.15]{1.15} \footnote[1]{Le lecteur qui ne
s'intéresserait qu'au cas noethérien est invité à se reporter à
\sisi{la remarque }{}\hyperref[XV.rem1.18]{1.18}.}.
% il y a deux énoncés numérotés 1.18 dans cet exposé.
Montrons d'abord qu'on peut supposer $Y$ strictement
local. Dans le cas abélien, $n\geq 1$, la suite spectrale de Leray
pour $g$ et $g^{\ast}(F)$, jointe à $R^{i}g_{\ast}(g^{\ast}(F))=0$
pour $0<i\leq n-1$, et $F\xrightarrow{\sim}g_{\ast}g^{\ast}(F)$,
exprimant l'hypothèse de $(n-1)$-acyclicité de $g$, implique qu'on a
une suite exacte canonique
$$
0\longrightarrow H^{n}(Y,F)\longrightarrow
H^{n}(X,g^{\ast}(F))\longrightarrow
H^{0}(Y,R^{n}g_{\ast}(g^{\ast}(F))){\quoi},
$$
qui montre qu'il suffit de prouver que l'image de $\xi$ dans le
troisième terme est nul. Ceci se vérifie fibre par fibre, et compte
tenu de \hyperrefext[VIII][VIII.thm5.2]{5.2} on est ramené au cas où $Y$ est strictement local.
Le cas $n=0$, $F$ étant un faisceau d'ensembles, se traite de la même
manière, car l'hypothèse sur $g$ implique que $F\rightarrow
g_{\ast}g^{\ast}(F)$ est injectif, et $\xi$ s'identifie à une section
du deuxième faisceau, dont il s'agit de montrer qu'elle provient d'une
section du premier: cela\pageoriginale se vérifie encore fibre par
fibre, et on applique \hyperrefext[VIII][VIII.rem5.3]{5.3}.
Le même argument essentiellement est
valable dans le cas non commutatif, avec $n=1$, en utilisant la suite
exacte \hyperrefext[XII][XII.prop3.2]{3.2}
$$
H^{1}(Y,F)\longrightarrow H^{1}(X,g^{\ast}(F))\longrightarrow
H^{0}(Y,R^{1}g_{\ast}(g^{\ast}(F))){\quoi},
$$
et utilisant \hyperrefext[VIII][VIII.rem5.3]{5.3}.

Nous supposons donc $Y$ strictement local, donc $X$ quasi-compact et
quasi-séparé. Le morphisme $g:X\rightarrow Y$ satisfait aux conditions
de \hyperrefext[XII][XII.cor6.10]{6.10}, ce qui nous ramène au cas où pour toute partie fermée
$Y_{1}$ de $Y$ distincte de $Y$, la restriction de $\xi$ à
$X_{1}=\fprod{X}{Y_1}{Y}$ est dans l'image de
$H^{n}(Y_1,F|Y_1)\rightarrow H^{n}(X_1,g^{\ast}(F)|X_1)$, et à
vérifier dans ce cas qu'on peut trouver un morphisme fini $f:Y' \rightarrow Y$,
dont l'image dans $Y$ contient un ouvert non vide de $Y$, et tel que
l'image inverse $\xi'$ de $\xi$ sur $X'=\fprod{X}{Y'}{Y}$ soit
contenue dans l'image de $H^{n}(Y',F')\rightarrow
H^{n}(X',{g'}^{\ast}(F'))$. On notera en effet que, compte tenu de
\hyperrefext[VIII][VIII.prop5.5]{5.5}, \hyperrefext[VIII][VIII.rem5.8]{5.8}, la condition $2^{\circ}$) de \loccit est
automatiquement satisfaite pour ce morphisme $f:Y'\rightarrow Y$.

Pour trouver un tel $f$, nous pouvons supposer $Y$ réduit, et nous
considérons un point maximal $y$ de $Y$. Soit $\bar{y}$ le point
géométrique correspondant à une clôture séparable de $k(y)$, de sorte
que $\bar{y}$ est limite projective filtrante d'ouverts $U_i$ de
schémas finis intègres $Y_i$ sur $Y$, dont chacun est tel que son
image dans $Y$ contient un voisinage de $y$. Utilisant l'hypothèse sur
$\xi$, on voit qu'il existe un indice $i$ tel que l'image inverse de
$\xi$ sur $\fprod{X}{U_i}{Y}$ est nulle si $n\geq 1$, \resp est dans
l'image de $H^{0}(U_i,F_{U_i})$ si $n=0$.

Nous\pageoriginale prendrons $Y'$ égal à $Y_i$, et notons qu'on peut
supposer $U'=U_i$ égal à l'image inverse d'un voisinage ouvert affine
$U$ de $y$ dans $Y$, de sorte qu'on voit que l'image inverse de $\xi$
sur $\fprod{X}{Z'}{Y}$ (où $Z'=Y'-U'$ est l'image inverse de $Z=X-U$
dans $Y'$) est dans l'image de $H^{n}(Z',F_{Z'})$. Quitte à remplacer
$Y$ par $Y'$, nous voyons donc que nous sommes réduits au cas suivant:
il existe un ouvert affine $U$ de $Y$, de complémentaire $Z$, tel que
les restrictions de $\xi$ à $X_U$ et à $X_Z$ appartiennent
respectivement à l'image de $H^{n}(U,F_U)$ et à l'image de
$H^{n}(Z,F_Z)$; de plus, si $n\geq 1$, on peut supposer ces
restrictions \emph{nulles}.

Lorsque $n\geq 1$, nous allons montrer que sous ces conditions, on a
$\xi=0$. Pour simplifier les notations, nous écrirons $Y'$ au lieu de
$X$, et nous désignons par $U'$, $Z'$ les images inverses de $U$, $Z$
dans $Y'$, par $F'$ l'image inverse $g^{\ast}(F)$. Notons
$i:U\rightarrow Y$, $j:Z\rightarrow Y$, $i':U'\rightarrow Y'$,
$j':Z'\rightarrow Y'$ les inclusions,
$$
\xymatrix@C=2cm@R=1.5cm{
U'  \ar[r]^{i'} \ar[d]^{g_U}& Y' = X  \ar[d]^{g} & Z'  \ar[l]_{j'} \ar[d]^{g_Z}\\
U   \ar[r]^{i} & Y & Z \ar[l]_{j},
}
$$
et considérons la suite exacte de faisceaux
$$
0\longrightarrow F'_{U',Y'}\longrightarrow F'\longrightarrow
F'_{Z',Y'}\longrightarrow 0{\quoi}.
$$
% référence pour cette notation ?
Comme l'image de $\xi$ dans $H^{n}(Y',F'_{Z',Y'})=H^{n}(Z',F'_{Z'})$
est nulle, $\xi$ provient d'un élément de $H^{n}(Y',F'_{U',Y'})$, et
la restriction de ce dernier sur $U'$ est égale à celle de $\xi$, donc
est nulle. Comme $F'_{U',Y'}$ est isomorphe à l'image inverse de
$F_{U,Y}$, nous sommes ainsi ramenés\pageoriginale au cas d'un
faisceau de la forme $F_{U,Y}$, donc nous pouvons supposer $F_Z=0$ (ce
qui tient compte de l'hypothèse $\xi|Z'=0$). Notons que si
$F\rightarrow G$ est un monomorphisme de faisceaux abéliens de
$\LL$-torsion (\resp de faisceaux en groupes), il suffit de
prouver que l'image $\eta$ de $\xi$ dans $H^{n}(Y', G')$ est dans
l'image de $H^{n}(Y,G)$ \ie est nulle \hyperrefext[XII][XII.cor6.6]{6.6}. Cela nous permet de
remplacer $F$ par $i_{\ast}(F_U)$, compte tenu que l'hypothèse $F_Z=0$
implique que l'homomorphisme canonique $F\rightarrow i_{\ast}(F_U)$
est injectif. Or comme $g$ est localement $(n-1)$-acyclique pour
$\LL$, et a fortiori $0$-acyclique pour $\LL$ (car
$n\geq 1$), il s'ensuit que l'image inverse de $i_{\ast}(F_U)$ sur
$Y'$ s'identifie à $i'_{\ast}(F'_{U'})$, de sorte qu'on est ramené à
montrer que l'image $\eta$ de $\xi$ dans
$H^{n}(Y',i'_{\ast}(F'_{U'}))$ est nulle. Or la restriction de $\eta$
à $U'$ est nulle, et il suffit donc de vérifier que l'homomorphisme
\begin{equation*}
H^{n}(Y',i'_{\ast}(F'_{U'}))\longrightarrow H^{n}(U',F'_{U'})\tag{$\ast$}
\end{equation*}
est un monomorphisme. Dans le cas abélien, on utilise la suite
spectrale de Leray pour $i'$ et $F'_{U'}$:
$$
E^{p,q}_2=H^p(Y',R^qi'_{\ast}(F'_{U'}))\Longrightarrow
H^{\ast}(U',F'_{U'}){\quoi},
$$
et les relations
\begin{equation*}
E^{p,q}_2=0\; \mbox{\rm pour } 0<p\leq n-1\quav q\leq n-1{\quoi}.\tag{${\ast}{\ast}$}
\end{equation*}
Pour vérifier ces dernières, on note que l'hypothèse de
$(n-1)$-acyclicité locale de $g$ pour $\LL$ implique que l'on a
des isomorphismes
$$
R^{q}i'_{\ast}(F'_{U'})\simeq g^{\ast}(R^{q}i_{\ast}(F_U))\;\mbox{\rm pour } q\leq n-1{\quoi},
$$
et\pageoriginale compte tenu du fait que $g$ est $(n-1)$-acyclique
pour $\LL$, donc que $H^{p}(Y,G)\xrightarrow{\sim}H^{p}(Y',G')$
pour $p\leq n-1$ et tout faisceau abélien de $\LL$-torsion $G$
sur $Y$, on trouve $({\ast}{\ast})$ puisque $H^{p}(Y,G)=0$ pour
$p>0$, et on trouve de même
\begin{equation*}
E^{0,q}_2\simeq H^{0}(Y,R{}^qi_{\ast}(F_U))\;\mbox{\rm
pour}\quad q\leq n-1{\quoi}.\tag{${\ast}{\ast}{\ast}$}
\end{equation*}
La suite spectrale de Leray considérée donne alors naissance à une
suite exacte, qui forme la deuxième ligne du diagramme commutatif
suivant:
$$
\xymatrix{
**[r]H^{0}(Y,R^{n-1} i_{\ast}(F_U))  \ar[r] \ar[d]^{\text{\rotatebox{90}{$\sim$}}}&
**[r]H^n(Y,i_{\ast} (F_U)) = 0 \ar[d]\\
**[r]H^{0}(Y', R^{n-1} i'_{\ast}(F'_{U'}))  \ar[r] &
**[r]H^{n}(Y',i'_{\ast}(F'_{U'}))  \ar[r] & **[r]H^{n}(U',F'_{U'}),
}
$$
où la première flèche verticale est bijective en vertu de
$({\ast}{\ast}{\ast})$. On lit sur le diagramme que $(\ast)$ est
injectif, ce qui achève la démonstration dans ce cas. Dans le cas non
abélien, $n=1$, on utilise le même diagramme avec $n=1$:
$$
\xymatrix{
**[r]H^{0}(Y, i_{\ast}(F_U))  \ar[r] \ar[d]^{\text{\rotatebox{90}{$\sim$}}}&
**[r]H^1(Y,i_{\ast} (F_U)) = 0 \ar[d]\\
**[r]H^{0}(Y', i'_{\ast}(F'_{U'}))  \ar[r] &
**[r]H^{1}(Y',i'_{\ast}(F'_{U'}))  \ar[r] & **[r]H^{1}(U',F'_{U'}),&
}
$$
\hyperrefext[XII][XII.prop3.2]{3.2}, où la première flèche verticale est encore bijective grâce à
l'hypothèse de $0$-acyclicité locale et globale faite sur $g$, et on
conclut encore l'injectivité de $(\ast)$ dans ce cas.

Traitons enfin le cas $n=0$. Remarquons que la donnée d'une section
$\xi$ de $F'$ revient, en vertu de \sisi{\qquad}{\qquad},
% référence manquante
à la donnée d'une section $\xi_{U'}$ de $F'_{U'}$ et d'une section $\xi_{Z'}$ de $F'_{Z'}$,
telles que l'image de $\xi_{Z'}$ par l'homomorphisme naturel
$$
F'_{Z'}\longrightarrow {j'}^{\ast}(i'_{\ast}(F'_{U'}))
$$
soit\pageoriginale la section du deuxième membre induite par
$\xi_{U'}$. Par hypothèse, il existe des sections $\eta_Z$ et $\eta_U$
de $F_Z$ et $F_U$ respectivement, qui induisent $\xi_{Z'}$ et
$\xi_{U'}$. Tout revient à voir qu'elles satisfont à la condition de
compatibilité analogue, exprimant que ce sont les restrictions d'une
même section $\eta$ de $F$. Il faut vérifier l'égalité de deux
sections de $j^{\ast}(i_{\ast}(F_U))$, et pour ceci, comme
$Z'\rightarrow Z$ est surjectif ($g$ étant surjectif puisque
$(-1)$-acyclique par hypothèse), il suffit de vérifier que les deux
sections correspondantes de
$g_Z{}^{\ast}(j^{\ast}(i_{\ast}(F_U)))=j'{}^{\ast}(g^{\ast}(i_{\ast}(F_U)))$
sont égales. Or le faisceau envisagé, grâce à l'homomorphisme de
changement de base
\begin{equation*}
g^{\ast}(i_{\ast}(F_U))\longrightarrow
i'_{\ast}(g_U{}^{\ast}(F_U))=i'_{\ast}(F'_{U'}){\quoi},\tag{$\ast$}
\end{equation*}
s'envoie dans ${j'}^{\ast}(i'_{\ast}(F'_{U'}))$, et ce dernier
homomorphisme est injectif, car il en est ainsi de $(\ast)$ grâce à
l'hypothèse que $g$ est localement $(-1)$-acyclique. Il suffit donc de
vérifier que les deux sections de ${j'}^{\ast}(i'_{\ast}(F'_{U'}))$,
images des deux sections précédentes, sont égales. Or on constate
aussitôt que ces sections ne sont autres que celles déjà envisagées
plus haut, qui sont égales par hypothèse. Cela achève la
démonstration.

\begin{enonce*}{Remarque 1.18}\etiquette[1.18]{XV.rem1.18}
On peut donner un raffinement de \hyperref[XV.thm1.15]{1.15}, qui fournit une démonstration
nettement plus simple de \hyperref[XV.thm1.15]{1.15} lorsque $Y$ est supposé localement
noethérien. Pour ceci, en plus des hypothèses préliminaires sur $g$,
$F$, supposons qu'il existe un nombre fini de points géométriques
$f_i:\bar{y}_i\rightarrow Y$ de $Y$ tels que le morphisme 
$$
F\longrightarrow \prod_if_{i\ast}(f_i{}^{\ast}(F))
$$
soit\pageoriginale injectif (hypothèse automatiquement vérifiée si $Y$
est noethérien et si $F$ est constructible \hyperrefext[IX][IX.prop2.14]{2.14}). Je dis que sous
ces conditions, le critère \hyperref[XV.thm1.15]{1.15} est valable en se bornant aux seuls
points géométriques $\bar{y}_i$. En effet, en vertu de \hyperrefext[XII][XII.cor6.6]{6.6} on peut
supposer que $F$ est lui-même de la forme $f_{\ast}(M_{\bar{y}})$, où
$M$ est un ensemble \resp un groupe \resp un groupe abélien de
$\LL$-torsion, et $f:\bar{y}\rightarrow Y$ un point
géométrique. Soit $f':X'=X_{\bar{y}}\rightarrow X$ le morphisme
canonique. Supposant, pour fixer les idées, $n\geq 1$, l'hypothèse
locale faite sur $g$ implique qu'on a
$$
g^{\ast}(f_{\ast}(M_{\bar{y}}))\simeq f'_{\ast}(M_{X'})\sisi{,\;}{\quav}
R^{i}f'_{\ast}(M_{X'})=0\quad\mbox{\rm pour}\quad 1\leq i\leq
n-1{\quoi},
$$
la deuxième relation provenant du fait que le premier membre est
isomorphe à\break $g^{\ast}(R^{i}f_{\ast}(M_{\bar{y}}))$, qui est nul
puisque manifestement $R^{i}f_{\ast}(M_{\bar{y}})=0$ pour $i\geq
1$. De ces relations, on conclut que
$$
H^{n}(X,g^{\ast}(F))\simeq H^{n}(X,f'_{\ast}(M_{X'}))\longrightarrow
H^{n}(X',M_{X'})
$$
est injectif, en utilisant la suite spectrale de Leray pour $f'$, ce
qui prouve bien qu'un élément $\xi$ de $H^{n}(X,g^{\ast}F)$ induisant un
élément nul de $H^{n}(X_{\bar{y}},F)$ est nul. Pour déduire de ceci
une démonstration simplifiée de \hyperref[XV.thm1.15]{1.15} dans le cas où $Y$ est localement
noethérien, on note qu'on peut se borner au cas $Y$ noethérien, et un
passage à la limite facile, utilisant \hyperrefext[IX][IX.cor2.7.2]{2.7.2} et \hyperrefext[VII][VII.sec5]{5},
permet de se ramener au cas où $F$ est de plus constructible, justiciable de
l'énoncé qu'on vient de prouver.
\end{enonce*}

\section{Acyclicité locale d'un morphisme lisse}\etiquette[2]{XV.sec2}\pageoriginale

Le théorème est le suivant~:

\begin{enonce*}{Théorème 2.1}\etiquette[2.1]{XV.thm2.1}
Soient $g:X\rightarrow Y$ un morphisme lisse et soit $\LL\subset
\PP$ l'ensemble complémentaire de l'ensemble des
caractéristiques résiduelles de $X$. Alors $g$ est localement
acyclique et localement $1$-asphérique pour $\LL$.
\end{enonce*}

Notons d'abord le corollaire:

\begin{enonce*}{Corollaire 2.2}\etiquette[2.2]{XV.cor2.2}
Soit $g:\EE^m_Y\rightarrow Y$ le morphisme structural de
l'espace affine. Alors $g$ est acyclique et $1$-asphérique pour
l'ensemble $\LL$ complémentaire de l'ensemble des
caractéristiques résiduelles de $Y$.
\end{enonce*}

En effet, par récurrence sur $m$ et \hyperref[XV.prop1.9]{1.9} (i) on se ramène au cas
$m=1$. D'après \hyperref[XV.thm2.1]{2.1} on sait que $g$ est localement $n$-acyclique et
localement $1$-asphérique pour $\LL$, donc que la condition
\emph{locale} de \hyperref[XV.cor1.16]{1.16} est satisfaite. Il suffit donc de démontrer que
les fibres géométriques de $\EE^{1}_Y$ sont acycliques et
$1$-asphériques pour $\LL$, c'est-à-dire que $\Spec k[t]$ est
acyclique et $1$-asphérique pour $\LL$ si $k$ est un corps
séparablement clos de caractéristique $p\not\in\LL$, ce qui
résulte de \hyperrefext[IX][IX.thm4.6]{4.6} pour l'assertion d'acyclicité, et pour la
$1$-asphéricité est bien connu. Plus généralement, soit $X'$ un
revêtement normal irréductible galoisien de $X=\PP^{1}_k$,
qui\pageoriginale soit non ramifié en dehors du point $\infty$ et dont
la ramification en ce point soit modérée. Prouvons que le degré $n$ du
revêtement est égal à $1$. Or on a la formule de Hurwitz:
$$
2(g'-1)=2n(g-1)+\deg \mathfrak{d}_{X'/X}
$$
où $g$, $g'$ sont le genre de $X$, $X'$, et $\mathfrak{d}_{X'/X}$ le diviseur
différente. On a $g=0$, et l'hypothèse de ramification modérée
implique $\deg \mathfrak{d}_{X'/X}\leq n-1$, d'où $2(g'-1)\leq -(n+1)$ donc
$n\leq 2(1-g')-1$, et comme $g'\geq 0$, on conclut $n\leq 1$ donc
$n=1$, \cqfd % CQFD à droite après une virgule est peu esthétique

\medskip
\noindent
\emph{Démonstration de} \hyperref[XV.thm2.1]{2.1}. L'assertion est locale sur $X$ et $Y$
pour la topologie étale, et on peut donc supposer \sisi{(}{}\SGA 1 II
1.1\sisi{)}{} que $g$ est le morphisme structural de l'espace affine $\EE^m_Y$,
et que $Y$ est affine. En appliquant la transitivité de l'acyclicité
locale \hyperref[XV.prop1.13]{1.13}.~(i) et le fait que
$\EE^m_Y=\EE^{m-1}_{\EE^1_Y}$ on se réduit
immédiatement au cas $m=1$, c'est-à-dire, au cas
$X=\Spec \mathcal{O}_Y[t]$. % Spec souligné dans l'original

Appliquons \hyperref[XV.cor1.18]{1.18}. On est ramené à démontrer que les fibres géométriques
du morphisme de schémas strictement locaux
$\widetilde{g}:\widetilde{X}\rightarrow \widetilde{Y}$, déduit d'un
couple de points géométriques $\bar{x}$, $\bar{y}$, où $\bar{x}$ est
fermé dans la fibre $X_{\bar{y}}$, sont acycliques pour $\LL$
et $1$-asphériques pour $\LL$. On peut évidemment supposer
$Y=\widetilde{Y}$. Soit $Y'\rightarrow Y$ un morphisme fini, local, et
surjectif, de sorte que $Y'$ est strictement local, et soient
$\bar{y}'$ le point de $Y'$ au-dessus de $\bar{y}$, $\bar{x}'$ le
point de $X'$ au-dessus de $\bar{x}$ et de $\bar{y}'$. Puisque
$X'\rightarrow X$ est fini, le localisé strict de $X'$ au point
$\bar{x}'$ n'est autre que
$\widetilde{X}'=\fprod{\widetilde{X}}{Y'}{Y}$. Par suite, toute fibre
géométrique de $\widetilde{X}'/Y'$ est déduite d'une fibre géométrique
de $\widetilde{X}/Y$ par extension algébrique
(nécessairement\pageoriginale radicielle) du corps de base, de plus il
y a au moins une fibre géométrique de $\widetilde{X}'/Y'$ au-dessus de
chaque fibre géométrique de $\widetilde{X}/Y$. On peut donc \hyperrefext[VIII][VIII.thm1.1]{1.1}
remplacer $Y$ par $Y'$, et par choix convenable de $Y'$, on se ramène
ainsi au cas où $k(\bar{x})=k(\bar{y})$. Alors $\bar{x}$ est un point
rationnel de la fibre $X_{\bar{y}}=\Spec k(\bar{y})[t]$, disons le
point $t=a^{0}$. En relevant $a^{0}$ en une section $a$ de
$\mathcal{O}_{\sisi{Y}{\widetilde{Y}}}$ et en faisant la « translation »  $t_1=t-a$ on
est ramené au cas où le point $\bar{x}$ est le point $\{t=0\}$ de
$\Spec k(\bar{y})[t]$. Posons $A=\Gamma
(\sisi{Y}{\widetilde{Y}},\mathcal{O}_{\sisi{Y}{\widetilde{Y}}})$, de sorte que $A$ est un anneau
strictement local.

\begin{enonce*}[remark]{Notation 2.3}\etiquette[2.3]{XV.not2.3}
Soit $A$ un anneau hensélien local. On dénote par $A\{t\}$ le
hensélisé de l'anneau $A[t]$ en l'idéal premier engendré par $\rad
A$ et $t$.
\end{enonce*}

\begin{enonce*}{Corollaire 2.4}\etiquette[2.4]{XV.cor2.4}
Soient $A$ hensélien et $A'$ une $A$-algèbre finie. Alors
$$
A'\{t\}\simeq A'\otimes_A A\{t\}{\quoi}.
$$
\end{enonce*}

Cela résulte immédiatement de \hyperrefext[VIII][VIII.subsec4.1]{4.1}.
Notons que si $A$ est
strictement local, il en est de même de $A\{t\}$.

La démonstration du théorème est ramenée au lemme suivant:

\begin{enonce*}{Lemme 2.5}\etiquette[2.5]{XV.lem2.5}
Soit $A$ un anneau hensélien de caractéristique résiduelle $p$. Alors
chaque fibre géométrique de $\Spec A\{t\}$ sur $\Spec A$ est acyclique
et $1$-asphérique pour $\LL=\PP-\{p\}$.
\end{enonce*}

Notons que puisque $A\{t\}$ est le hensélisé de $A[t]$, une fibre
$\Spec A\{t\}\otimes_A K$, $K$ un corps résiduel de $A$, est limite
projective\pageoriginale de schémas affines et de type fini $X_i$ sur
la « droite »  $\Spec K[t]$, c'est-à-dire, limite de courbes
algébriques affines. Donc chaque fibre géométrique de $\Spec A\{t\}$
est limite de courbes affines $(X_i)_{\bar{K}}$ au-dessus du corps
séparablement clos $\bar{K}$, et il s'ensuit de \hyperrefext[IX][IX.cor5.7]{5.7} et \hyperrefext[VII][VII.thm5.7]{5.7}
que la dimension cohomologique des fibres géométriques est $\leq 1$. Il
suffit donc de vérifier que les fibres géométriques sont
$1$-\emph{asphériques} pour $\LL$, c'est-à-dire, qu'une telle
fibre géométrique $\bar{Z}$ est connexe, non vide, et que
$H^{1}(\bar{Z},G)=0$ pour chaque $\LL$-groupe fini $G$.

Écrivons $A=\varinjlim B_{\alpha}$ où les $B_{\alpha}$ sont des
anneaux de type fini sur $\ZZ$. Le morphisme
$B_{\alpha}\rightarrow A$ se factorise par le hensélisé $A_{\alpha}$
de $B_{\alpha}$ en l'idéal premier induit par $\rad A$, d'où
$A=\varinjlim A_{\alpha}$. On constate aisément qu'on a aussi
$A\{t\}=\varinjlim A_{\alpha}\{t\}$. De plus, soient $\bar{y}$ un
point géométrique de $\Spec A$, $X_{\bar{y}}$ la fibre géométrique de
$S=\Spec A\{t\}$ sur $\Spec A$ en $\bar{y}$, $\bar{y}_{\alpha}$ le
point de $\Spec A_{\alpha}$ induit par $\bar{y}$, et
$X_{\alpha}\,_{\bar{y}_{\alpha}}$ % \, : à changer ?
la fibre géométrique de $\Spec A_{\alpha}\{t\}$ sur $\Spec A_{\alpha}$ en~$\bar{y}_{\alpha}$. On a
$X_{\bar{y}}=\varprojlim {X_{\alpha}}\,_{\bar{y}_{\alpha}}$. On est donc
réduit pour la démonstration de \hyperref[XV.lem2.5]{2.5} par \hyperref[XV.prop1.9]{1.9} (ii) (appliqué aux
morphismes $X_{\alpha}\,_{\bar{y}_{\alpha}}\rightarrow \bar{y}_{\alpha}$)
au cas $A=A_{\alpha}$, c'est-à-dire, où $A$ est le hensélisé d'une
algèbre de type fini sur $\ZZ$ en un idéal premier donné. Nous
le supposons désormais. Rappelons qu'une telle algèbre est un anneau
\emph{excellent} \sisi{(}{}\EGA IV 7.8.3 (iii), 7.8.6 (i)\sisi{)}{}, donc que $A_{\alpha}$
est également excellent (comme on voit facilement sur la définition,
\cf \EGA IV 18.7.6).

Vérifions la $0$-acyclicité. (Pour un résultat plus général, \cf
Appendice.) Puisque $A\{t\}/A$ admet la section « $t=0$ », une
fibre géométrique\pageoriginale $\bar{Z}$ est non-vide. Pour voir
qu'elle est connexe, on est réduit, en remplaçant $A$ par un
quotient \sisi{(}{}\hyperref[XV.cor2.4]{2.4}\sisi{)}{}, au cas où $A$ est intègre, et où $\bar{Z}$ est la fibre
géométrique au-dessus du point générique de $\Spec A$. Il suffit de
prouver que pour toute extension finie séparable $K'$ de corps de
fractions $K$ de $A$, le schéma $Z_{K'}$ ($Z=\Spec A\{t\}\otimes
{}_AK$) est connexe. Soient $A'$ la normalisée de $A$ dans $K'$ (qui
est une $A$-algèbre finie \sisi{(}{}\EGA IV 7.8.2\sisi{)}{})
et $Z'$ la fibre générique de $\Spec A'\{t\}$ sur $A'$. D'après \hyperref[XV.cor2.4]{2.4},
on a $A'\{t\}\simeq A'\otimes_A A\{t\}$, d'où $Z_{K'}\simeq Z'$. Remplaçons $A$ par $A'$. Il suffit
alors de démontrer que $Z$ est intègre, et il suffit de démontrer que
$A\{t\}$ est intègre. Mais $A$ étant normal, $A[t]$ l'est, donc
$A\{t\}$ l'est aussi \sisi{(}{}\SGA 1 I 9.5.~(i)\sisi{)}{}, d'où le résultat.

La démonstration de \hyperref[XV.lem2.5]{2.5}, et par suite de \hyperref[XV.thm2.1]{2.1},
est ainsi réduite au lemme suivant, qui sera démontré dans le prochain numéro:

\begin{enonce*}{Lemme 2.6}\etiquette[2.6]{XV.lem2.6}
Soit $A$ un anneau hensélien et excellent, et soit $\bar{Z}$ une fibre
géométrique de $\Spec A\{t\}/\Spec A$. Alors chaque revêtement
principal galoisien $\bar{Z}'$ de $\bar{Z}$, de groupe $G$ d'ordre
premier à la caractéristique résiduelle $p$ de $A$, est trivial.
\end{enonce*}

\section{Démonstration du lemme principal}\etiquette[3]{XV.sec3}\pageoriginale

En remplaçant $A$ par une algèbre finie $A'$ (\cf \hyperref[XV.cor2.4]{2.4}), on réduit
l'assertion \hyperref[XV.lem2.6]{2.6} au cas où $A$ est intègre et où $\bar{Z}$ est la fibre
géométrique générique de $\Spec A\{t\}$ sur $\Spec A$. Soient $K$ le
corps des fractions de $A$ et $Z=\Spec A\{t\}\otimes_A K$. Si
$\bar{Z}'/\bar{Z}$ est un revêtement principal galoisien, on peut le
descendre à un $Z_{K'}$, où $K'$ est une extension finie séparable
convenable de $K$. En remplaçant \hyperref[XV.cor2.4]{2.4} $A$ par son normalisé dans
$K'$, qui est une $A$-algèbre finie \sisi{(}{}\EGA IV 7.8.3 (vi)\sisi{)}{}, on est ramené à
démontrer le lemme sous la forme suivante:

\begin{enonce*}{Lemme 3.1}\etiquette[3.1]{XV.lem3.1}
Avec les notations de \hyperref[XV.lem2.6]{2.6}, supposons $A$ normal, et soit $Z'/Z$ un
revêtement étale principal galoisien, de groupe $G$ d'ordre premier à
la caractéristique résiduelle $p$ de $A$. Alors $Z'$ est induit par
une extension galoisienne $K'$ de $K$, \ie \emph{(\SGA  1 IX 6.2)} le
revêtement $\bar{Z}'$ de $\bar{Z}$ déduit de $Z'$ est trivial.
\end{enonce*}

\noindent
\emph{Démonstration}. Soit $B'$ le normalisé de $A\{t\}$ dans l'anneau
des fonctions rationnelles de $Z'$, et considérons les idéaux premiers
$P$ de hauteur $1$ de $A\{t\}$ tels que l'algèbre $B'$ sur $A\{t\}$
soit ramifiée en $P$. Puisque $Z'/Z$ est étale, $P$ n'est pas sur le
point générique de $\Spec A$, \ie $\mathfrak{p}=P\cap A\neq 0$. Donc
$A\{t\}/\mathfrak{p}A\{t\}=(A/\mathfrak{p})\{t\}$ est intègre et de dimension au plus
égale à $\dim A$. Puisque $A\{t\}$ est excellent \sisi{(}{}\EGA IV 7.8.6 (i) et
18.7.6\sisi{)}{} donc caténaire, on a $\dim \Spec A\{t\}/P=(\dim A+1)-1=\dim
A\geq \dim A\{t\}/\mathfrak{p}A\{t\}$, d'où\pageoriginale
$P=\mathfrak{p}A\{t\}$ puisque $P\supset \mathfrak{p}A\{t\}$.

Soient donc $\mathfrak{p}_1,\ldots,\mathfrak{p}_m$ les idéaux premiers de hauteur~$1$
de~$A$ tels que $B'/A\{t\}$ soit ramifié en les idéaux $\mathfrak{p}_iA\{t\}$,
soit $s$ le p.p.c.m. des ordres des groupes d'inertie \sisi{(}{}\SGA
1 V 2\sisi{)}{} $G_{i}\subset G$ pour $B'$ en les idéaux premiers au-dessus des
$\mathfrak{p}_{i}A\{t\}$, qui est un entier divisant l'ordre de $G$, donc
premier à $p$. Soit $f\in A$ une fonction qui s'annule avec ordre~$1$ en
chaque $\mathfrak{p}_i$ (un tel $f$ existe, comme il résulte par exemple de
Bourbaki, Alg. Comm., Chap. II, \S 3, cor. à prop. 17). En remplaçant
\sisi{(}{}\hyperref[XV.cor2.4]{2.4}\sisi{)}{} $A$ par son normalisé $A_i$ dans l'extension finie
$K_{1}=K[x]/(x^{s}-f)$ de $K$ et $B'$ par le normalisé de
$A_{1}\otimes_A B'$, on se réduit par le lemme d'Abhyankar
\sisi{(}{}\SGA 1 X 3.6\sisi{)}{}
au cas où $B'$ n'est ramifiée en aucun idéal premier de hauteur~$1$ de $A\{t\}$.

Alors l'algèbre $B'/A\{t\}$ est aussi étale au-dessus de la fibre de
$\Spec A\{t\}$ en tout point de $\Spec A$ de codimension~$1$. En effet,
soit $\mathfrak{p}$ un idéal premier de hauteur~$1$ de $A$. Puisque $A$ est
normal, $A_{\mathfrak{p}}$ est un anneau de valuation discrète, donc
régulier, donc $A\{t\}\otimes_A A_{\mathfrak{p}}$ est aussi régulier
---~en effet, c'est une limite de schémas étales au-dessus de $A_{\mathfrak{p}}[t]$.
Par hypothèse, l'algèbre $B'/A\{t\}$ est étale en chaque point
de codimension $\leq 1$ de $\Spec A\{t\}\otimes_A A_{\mathfrak{p}}$, donc
elle est étale partout par le théorème de pureté de Zariski-Nagata
\sisi{(}{}\SGA 1 X 3.1\sisi{)}{}; pour une démonstration, voir p. ex. \SGA 2, X 3.4.

En remplaçant \sisi{(}{}\hyperref[XV.cor2.4]{2.4}\sisi{)}{} $A$ par une $A$-algèbre finie convenable, on se
ramène au cas où \emph{de plus} l'algèbre $B'/tB'$ sur
$A\{t\}/tA\{t\}\simeq A$ est complètement décomposée au-dessus du
point générique de $\Spec A$. Donc il suffit de démontrer que sous ces
conditions, l'algèbre $B'$ sur $A\{t\}$\pageoriginale est
« triviale », et puisque $A\{t\}$ est hensélien, il suffit de démontrer que
$B'$ sur $A\{t\}$ est étale; en effet, il résulte du fait que
l'algèbre $B'/tB'$ splitte au-dessus du point générique de $\Spec A$
qu'elle splitte sur $\Spec(A)$ (\SGA 1 I 10.1\sisi{)}{}, donc que les extensions
résiduelles de $B'$ sur l'idéal maximal de $A\{t\}$ sont triviales, et
on applique (\EGA IV 18.5.14). On est donc ramené au lemme suivant:

\begin{enonce*}{Lemme 3.2}\etiquette[3.2]{XV.lem3.2}
Soit $B$ un anneau local excellent et normal, et soit $0\neq t\in \rad
B$ tel que $B/tB$ soit normal. Soit $B'$ une $B$-algèbre finie intègre
normale, contenant $B$, telle que:
\begin{enumerate}
\renewcommand{\theenumi}{\roman{enumi}}
\renewcommand{\labelenumi}{(\theenumi)}
\item L'algèbre $B'/tB'$ sur $B/tB$ est complètement décomposée en les
points maximaux de $\Spec(B/tB)$.

\item L'algèbre $B'$ sur $B$ est étale en les points de codimension
$1$ de $\Spec(B/tB)$.
\end{enumerate}
Alors $B'$ est étale sur $B$.
\end{enonce*}

\noindent
\emph{Démonstration}. En vertu de (ii), l'assertion est triviale pour
$\dim B\leq 2$, c'est-à-dire pour $\dim B/tB\leq 1$. Supposons que
$B'/B$ ne soit pas étale, soit $x$ un point maximal de l'ensemble des
points de $\Spec B/tB$ au-dessus duquel $B'/B$ n'est pas étale, et
soit $C$ l'anneau local de $\Spec B$ en $x$. Alors $C$ est encore un
anneau excellent et normal (\EGA IV 7.8.2), on a $0\neq t\in\rad C$, et
(i) et (ii) sont vrais pour la $C$-algèbre $C'=C\otimes {}_BB'$. On
peut donc supposer $B=C$, c'est-à-dire, que $B'$ est étale sur $B$
au-dessus de chaque point de $\Spec B/tB$, sauf peut-être à
l'origine. Puisque $B/tB$ est normal\pageoriginale (en fait unibranche
suffirait), la condition (i) implique que l'algèbre $B'/tB'$ est
complètement décomposée en dehors de l'origine. De plus, on peut
supposer $\dim B\geq 3$. On est donc ramené au lemme suivant:

\begin{enonce*}{Lemme 3.3}\etiquette[3.3]{XV.lem3.3}
Soit $B$ un anneau local excellent et normal, $0\neq t\in \rad B$,
$\dim B\geq 3$. Soit $B'$ une $B$-algèbre finie normale intègre,
contenant $B$, telle que l'algèbre $B'/tB'$ de $B/tB$ soit
complètement décomposée en dehors de l'origine. Alors $B'$ est étale
sur $B$.
\end{enonce*}

\noindent
\emph{Démonstration}. Soient $X=\Spec B$, $X_1=X-\{x_{0}\}$, où
$x_{0}$ est le point fermé, et soit $i:U\rightarrow X$ l'ouvert
des $x\in X$ en lesquels l'algèbre $B'/B$ est étale. Alors $(X-U)\cap
V(t)=\{x_{0}\}$ est de dimension $0$. Puisque $V(t)$ est défini
par \emph{une} équation, il s'ensuit que $\dim(X-U)\leq 1$ (\EGA
$0_{\text{\rm IV}}$ 16.2.5), donc ($X$ étant caténaire) $x\in
X-U\Rightarrow \dim \mathcal{O}_{X,x}\geq \dim B-1\geq 2$, \ie
$\codim(X-U,X)\geq 2$.

Soit $X'=\Spec(B')$, $U'$ l'image inverse de $U$ dans $X'$, alors, $X$
étant universellement caténaire, la relation $\codim(X-U,X)\geq 2$
implique la relation $\codim(X'-U',X')\geq 2$ (\EGA IV 5.6.10), donc,
$X'$ étant normal, on a
$$
B'=\Gamma(X',\Oo_{X'})\xrightarrow{\sim}
\Gamma(U',\Oo_{U'}){\quoi}.
$$
Soit alors $\mathcal{B}'=\widetilde{B}'$ le faisceau de $\Oo_X$-algèbres qui
définit $X'$ sur $X$, $\mathcal{C}$ sa restriction à $U$, alors la formule
précédente équivaut à la formule
$$
\mathcal{B}'\xrightarrow{\sim}i_{\ast}(\mathcal{C}){\quoi},
$$
et\pageoriginale \hyperref[XV.lem3.3]{3.3} se ramène par suite à l'énoncé suivant (où il
n'est plus question de $B'$): Soient $B$, $t$ comme dans 3.3,
$X=\Spec(B)$, $Y=\Spec(B/tB)$, $U$ un voisinage ouvert de
$Y_1=Y-\{x_{0}\}$ dans $X_1=X-\{x_{0}\}$, $U'$ un revêtement
étale de $U$ dont la restriction à $Y_1$ soit complètement décomposée,
$\mathcal{C}$ le faisceau de $\Oo_{U}$-algèbres définissant $U'$, alors
($i:U\rightarrow X$ désignant l'immersion canonique) $i_{\ast}(\mathcal{C})$
est une Algèbre cohérente \emph{étale} sur $X$ (ou, ce qui revient au
même, grâce à la relation $\prof \Oo_{X,x}\geq 2$ pour $x\in X-U$
signalée plus haut et à \EGA IV 5.10.5, $U'$ se prolonge en un
revêtement \emph{étale} $X'$ de $X$).

Soit $\widehat{X}=\Spec\widehat{B}$, où $\widehat{B}$ est le
complété de $B$. Alors $\widehat{B}$ est normal (\EGA IV 7.8.3 (v))
et $\widehat{B}/B$ est fidèlement plat. Soit
$$
\xymatrix@C=1.5cm@R=1cm{
\hat{U}      \ar[r]^{\sisi{\hat{i}}{\hat{i}}} \ar[d]^{f_1}& \hat{X} \ar[d]^{f}\\
U            \ar[r]^{i} &  X
}
$$
le diagramme cartésien déduit de $\widehat{X}\rightarrow X$. Puisque
$f$ est fidèlement plat, le foncteur $i_{\ast}$ commute au changement
de base~$f$ \sisi{(}{}\EGA IV 2.3.1\sisi{)}{}, et on est ramené \sisi{(}{}\EGA IV 17.7.1\sisi{)}{}
à prouver que $\hat{i}_{\ast}(\widehat{\mathcal C})$ % C^ non défini dans le texte
est une algèbre étale sur $\widehat{X}$. On peut donc remplacer $B$ par $\widehat{B}$,
c'est-à-dire, on peut supposer $B$ \emph{complet}. Dans ce cas, nous
allons prouver que $U'$ est un revêtement complètement décomposé de
$U$ (ce qui impliquera évidemment qu'il se prolonge en un revêtement
étale de $X$, et achèvera la démonstration). Or cette assertion est
maintenant conséquence immédiate de la « théorie de
Lefschetz locale » \sisi{(}{}\SGA\pageoriginale 2 X 2.1 (i) et 2.5\sisi{)}{}, comme on constate
immédiatement, compte 
tenu encore une fois de la relation $\prof \Oo_{X,x}\geq 2$ pour
$\dim \{x\}=1$.

\begin{enonce*}[remark]{Remarque 3.4}\etiquette[3.4]{XV.rem3.4}
On peut prouver \hyperref[XV.lem3.2]{3.2} et \hyperref[XV.lem3.3]{3.3} sous des conditions 
plus générales, en
suivant essentiellement la même démonstration: au lieu de supposer $B$
excellent et normal, il suffit de supposer que le complété
$\widehat{B}$ est normal (condition qui est stable par localisation,
grâce aux résultats de \EGA IV 7), ou seulement que le complété de $B$
pour la topologie $(tB)$-adique soit normal, et quotient d'un anneau
régulier.
\end{enonce*}

\section{Appendice: Un critère de $0$-acyclicité locale}\etiquette[4]{XV.sec4}\pageoriginale 

Les résultats du présent numéro ne seront plus utilisés dans la suite
du séminaire.

\begin{enonce*}{Théorème 4.1}\etiquette[4.1]{XV.thm4.1}
Soit $g:X\rightarrow Y$ un morphisme de schémas. On suppose $g$ plat
et à fibres séparables, et $X$ et $Y$ localement noethériens
\sisi{(}{}\resp $g$ localement de présentation finie\sisi{)}{}.
Alors $g$ est universellement localement $0$-acyclique.
\end{enonce*}

On peut évidemment supposer $X$ et $Y$ affines, ce qui permet, dans le
cas respé, de se ramener au cas où $Y$ donc $X$ est noethérien, par la
méthode standard de \EGA IV 8, utilisant ici \EGA IV 9.7.7. On peut donc
se borner à prouver le premier énoncé. L'hypothèse étant stable par
changement de base de type fini, on est réduit, compte tenu du passage
à la limite \hyperref[XV.prop1.13]{1.13} (ii), à prouver que $g$ est localement $0$-acyclique,
et pour ceci on est ramené par \hyperref[XV.cor1.17]{1.17} à prouver que si on suppose de
plus $g$ strictement local, alors les fibres géométriques de $g$ sont
connexes. Or c'est ce que dit \EGA IV 18.9.8.

\begin{enonce*}{Corollaire 4.2}\etiquette[4.2]{XV.cor4.2}
Soit $g:X\rightarrow Y$ un morphisme de schémas, avec $Y$
discret. Alors $g$ est universellement localement $0$-acyclique.
\end{enonce*}

On peut supposer que $Y$ est le spectre d'un corps parfait $k$
\hyperrefext[VIII][VIII.thm1.1]{1.1}.
D'autre part, on peut supposer $X$ affine, donc limite projective
de schémas affines de type fini sur $k$, et compte tenu de \hyperref[XV.prop1.13]{1.13}~(ii)
on est ramené au cas où $X$ est de type fini sur~$k$.
Enfin,\pageoriginale on peut évidemment supposer $X$ réduit, donc
séparable sur $k$ puisque $k$ est parfait. La conclusion résulte alors
de \hyperref[XV.thm4.1]{4.1}.

\begin{enonce*}{Corollaire 4.3}\etiquette[4.3]{XV.cor4.3}
Supposons que $g$ soit surjectif, et satisfasse aux conditions de
\hyperref[XV.thm4.1]{4.1} ou de \hyperref[XV.cor4.2]{4.2}; dans le cas non respé de \hyperref[XV.thm4.1]{4.1} on suppose de plus
que $g$ est, soit quasi-compact et quasi-séparé, soit universellement
ouvert. Alors $g$ est un morphisme de descente effective universelle
pour la catégorie fibrée des faisceaux étales sur des préschémas
variables.
\end{enonce*}

Dans les cas envisagés, $g$ est universellement submersif et il
résulte de \hyperrefext[VIII][VIII.prop9.1]{9.1} que $g$ est un morphisme de descente universelle
pour la catégorie fibrée envisagée. Cela nous permet, pour la question
d'effectivité, de nous borner au cas $Y$ affine. On voit de plus, dans
chacun des cas envisagés, que $X$ se recouvre par un nombre
\emph{fini} d'ouverts affines $X_i$ dont les images recouvrent
$Y$. Remplaçant $X$ par la somme disjointe des $X_i$, on peut alors
supposer $X$ affine, a fortiori quasi-compact et quasi-séparé sur
$Y$. On conclut alors grâce au raisonnement de \hyperrefext[VIII][VIII.subsubsec9.4.1]{9.4.1}, en
utilisant le fait \hyperrefext[XVI][XVI.thm1.1]{1.1} que pour un morphisme quasi-compact et
quasi-séparé $g$, la formation de $g_{\ast}(F)$ ($F$ un faisceau étale
sur $X$) commute à tout changement de base $Y'\rightarrow Y$ qui est
localement $0$-acyclique.

\begin{enonce*}[remark]{Remarques 4.4}\etiquette[4.4]{XV.rems4.4}
\begin{enumerate}
\renewcommand{\theenumi}{\alph{enumi}}
\renewcommand{\labelenumi}{\theenumi)}
\item On obtient ainsi la démonstration (qui avait été laissée en
suspens) de \hyperrefext[VIII][VIII.thm9.4]{9.4} d), comme cas particulier de \hyperref[XV.cor4.3]{4.3}. % étiquette pour d) ?
Une démonstration différente plus facile, n'utilisant pas le résultat
assez délicat\pageoriginale de \EGA IV 18.9.8, s'obtiendrait en notant
que grâce à \hyperrefext[VIII][VIII.prop9.1]{9.1} on peut se borner au cas du morphisme
$\Spec(K)\rightarrow \Spec(k)$ induit par une extension de corps
$K/k$, or $\Spec(K)$ est universellement localement $0$-acyclique sur
$\Spec(k)$, car il est même universellement acyclique pour
$\LL=\PP-\{p\}$, $p=\car k$, comme on verra dans \hyperrefext[XVI][XVI.cor1.5]{1.5}.

\item On ignore si tout morphisme quasi-compact et quasi-séparé,
surjectif et localement $(-1)$-acyclique (ou universellement
localement $(-1)$-acyclique) est un morphisme de descente effective
pour la catégorie fibrée des faisceaux étales sur des préschémas
variables. Cet énoncé semble assez plausible, et améliorerait \hyperref[XV.cor4.3]{4.3} et
\hyperrefext[VIII][VIII.thm9.4]{9.4} c), d). % étiquette pour c) et d) ?
On le rapprochera de l'énoncé d'effectivité dans
\textsc{Murre}, Sém. Bourbaki ${\rm n^{\circ}}$~293, p.~17.

\item Il est plausible que sous les conditions de \hyperref[XV.cor4.2]{4.2}, $g$ est même
localement acyclique pour $\LL$ et localement $1$-asphérique pour $\LL$,
où $\LL$ est l'ensemble des nombres premiers distincts des
caractéristiques résiduelles de $Y$. On peut dans cette question
supposer évidemment $Y$ spectre d'un corps algébriquement clos $k$, et
on peut montrer (\SGA 1964/65) que la réponse est affirmative
lorsqu'on dispose de la résolution des singularités pour les schémas
de type fini sur $k$. Donc la réponse est affirmative lorsque $Y$ est
de caractéristique nulle, comme on voit en utilisant les résultats de
\textsc{Hironaka}~\footnote[1]{\Cf \SGA 1 XIII.}
\end{enumerate}
\end{enonce*}


%XII 4.1, 4.2 equation ref.
%XII 1.2, 1.6 no ref. at 501

\ifx\danslelivre\undefined
\end{document}
\fi
